% =====================================================================
%  CHƯƠNG 3 — CƠ SỞ LÝ THUYẾT VÀ CÔNG NGHỆ SỬ DỤNG
%
%  File chương theo cấu trúc template Overleaf (hướng Ứng dụng).
%  Cách dùng: đặt file này trong thư mục Chuong/ của template và thêm
%      % =====================================================================
%  CHƯƠNG 3 — CƠ SỞ LÝ THUYẾT VÀ CÔNG NGHỆ SỬ DỤNG
%
%  File chương theo cấu trúc template Overleaf (hướng Ứng dụng).
%  Cách dùng: đặt file này trong thư mục Chuong/ của template và thêm
%      % =====================================================================
%  CHƯƠNG 3 — CƠ SỞ LÝ THUYẾT VÀ CÔNG NGHỆ SỬ DỤNG
%
%  File chương theo cấu trúc template Overleaf (hướng Ứng dụng).
%  Cách dùng: đặt file này trong thư mục Chuong/ của template và thêm
%      % =====================================================================
%  CHƯƠNG 3 — CƠ SỞ LÝ THUYẾT VÀ CÔNG NGHỆ SỬ DỤNG
%
%  File chương theo cấu trúc template Overleaf (hướng Ứng dụng).
%  Cách dùng: đặt file này trong thư mục Chuong/ của template và thêm
%      \include{Chuong/Chuong3}
%  vào main.tex (sau Chương 2).
%
%  GHI CHÚ VỀ HÌNH VẼ: dùng macro \figph{<tên-file>}{<tỉ-lệ-rộng>}{<chú-thích>}{<nhãn>}
%  (đã khai báo ở Chương 2, khai báo lại bằng \providecommand để an toàn nếu
%  biên dịch độc lập). Macro tự chèn ảnh thật nếu tệp tồn tại trong images/,
%  ngược lại hiển thị khung giữ chỗ để báo cáo vẫn biên dịch được.
%
%  KHÔNG hardcode số chương/mục/bảng/hình/công thức — luôn dùng \label + \ref / \cite.
% =====================================================================

% --- Macro tiện ích (idempotent — không định nghĩa lại nếu đã có) ------
\providecommand{\code}[1]{\texttt{\detokenize{#1}}} % in định danh snake_case an toàn
\providecommand{\figph}[4]{%
  \begin{figure}[htbp]
    \centering
    \IfFileExists{images/#1}%
      {\includegraphics[width=#2\textwidth,height=0.82\textheight,keepaspectratio]{#1}}%
      {\setlength{\fboxsep}{10pt}\fbox{\begin{minipage}[c][3.6cm][c]{#2\textwidth}\centering\itshape Sơ đồ minh hoạ --- chèn tệp ảnh \code{images/#1} vào thư mục \code{images/}.\end{minipage}}}%
    \caption{#3}%
    \label{#4}%
  \end{figure}}

\chapter{CƠ SỞ LÝ THUYẾT VÀ CÔNG NGHỆ SỬ DỤNG}
\label{chuong:co-so-ly-thuyet}

Chương~\ref{chuong:phan-tich-thiet-ke} đã trình bày kết quả phân tích và
thiết kế hệ thống. Chương này hệ thống hoá những \textit{cơ sở lý thuyết}
và \textit{công nghệ} làm nền tảng cho các quyết định thiết kế và hiện thực
hoá đó, qua đó lý giải vì sao mỗi công nghệ được lựa chọn. Trước hết,
mục~\ref{sec:co-so-ly-thuyet} trình bày các cơ sở lý thuyết cốt lõi: ôn tập
ngắt quãng và thuật toán SM-2, khung trình độ CEFR, kiến trúc REST, ánh xạ
đối tượng--quan hệ và mô hình ngôn ngữ lớn. Tiếp theo,
mục~\ref{sec:tong-quan-cong-nghe} tổng quan ngăn xếp công nghệ và
mục~\ref{sec:kien-truc-tong-quan} phân tích các kiểu kiến trúc được áp dụng.
Các mục~\ref{sec:cong-nghe-csdl}--\ref{sec:cong-cu-phat-trien} đi sâu lần
lượt vào công nghệ cơ sở dữ liệu, frontend, backend, cơ chế xác thực và
phân quyền, công nghệ trí tuệ nhân tạo và bộ công cụ phát triển.

% =====================================================================
\section{Cơ sở lý thuyết}
\label{sec:co-so-ly-thuyet}

\subsection{Ôn tập ngắt quãng và đường cong lãng quên}
\label{subsec:srs-theory}

Cơ sở khoa học của động cơ học trong hệ thống là hiện tượng
\textit{đường cong lãng quên} (forgetting curve) do Ebbinghaus phát hiện:
khả năng ghi nhớ một thông tin mới suy giảm theo hàm mũ theo thời gian nếu
không được củng cố \cite{ebbinghaus1913memory}. Một cách mô hình hoá đơn
giản, mức ghi nhớ \( R \) còn lại sau khoảng thời gian \( t \) kể từ lần
học gần nhất được biểu diễn qua công thức~\eqref{eq:forgetting}:

\begin{equation}
  \label{eq:forgetting}
  R(t) = e^{-\,t/S},
\end{equation}

trong đó \( S \) là \textit{độ bền trí nhớ} (memory stability) --- càng ôn
tập đúng lúc, \( S \) càng lớn và thông tin được giữ lâu hơn. Hệ quả thực
tiễn là: thay vì học nhồi nhét, người học nên ôn lại mỗi từ tại những
\textit{thời điểm đến hạn} được giãn dần. Các nghiên cứu về hiệu ứng giãn
cách (spacing effect) khẳng định ôn tập ngắt quãng cải thiện rõ rệt khả
năng ghi nhớ dài hạn so với học dồn \cite{cepeda2006distributed}. Đây là
nền tảng lý thuyết cho \textit{hệ thống ôn tập ngắt quãng} (Spaced
Repetition System --- SRS) mà hệ thống hiện thực hoá.

\subsection{Thuật toán lập lịch ôn tập SM-2}
\label{subsec:sm2-theory}

Để quyết định \textit{khi nào} ôn lại mỗi từ, hệ thống áp dụng thuật toán
\textbf{SM-2} có nguồn gốc từ dự án SuperMemo \cite{wozniak1990optimization}.
Mỗi cặp (người dùng, từ) được mô tả bằng ba đại lượng: hệ số dễ nhớ
\textit{easiness factor} \( EF \) (khởi tạo \( 2{,}5 \)), số lần ôn đúng
liên tiếp \( n \) và khoảng cách ôn tập \( I \) (tính theo ngày). Sau mỗi
lượt ôn, người học (ở đây là kết quả trả lời trong phiên) được quy về một
\textit{điểm chất lượng} \( q \in \{0,1,\dots,5\} \).

Hệ số \( EF \) được cập nhật theo công thức~\eqref{eq:sm2-ef}, có chặn dưới
\( 1{,}3 \) để khoảng ôn tập không co lại quá nhanh:

\begin{equation}
  \label{eq:sm2-ef}
  EF' = EF + \big(0{,}1 - (5-q)\,(0{,}08 + (5-q)\cdot 0{,}02)\big),
  \qquad EF' \geq 1{,}3 .
\end{equation}

Khoảng cách ôn tập kế tiếp được tính theo công thức~\eqref{eq:sm2-interval}:

\begin{equation}
  \label{eq:sm2-interval}
  I(n) =
  \begin{cases}
    1, & n = 1,\\[2pt]
    6, & n = 2,\\[2pt]
    \big\lceil I(n-1)\times EF' \big\rceil, & n > 2 .
  \end{cases}
\end{equation}

Nếu người học trả lời sai (\( q < 3 \)), bộ đếm \( n \) được đặt lại về 0 và
từ phải học lại từ đầu; ngược lại \( n \) tăng và khoảng ôn tập giãn dần
theo \( EF' \). Trong hệ thống, SM-2 được \textbf{mở rộng bằng các bước học}
(learning steps): một từ mới phải vượt qua một chuỗi câu hỏi từ dễ đến khó
trong phiên trước khi được giao cho SM-2 lập lịch dài hạn; điểm chất lượng
\( q \) được suy ra từ độ chính xác ở bước cuối của mỗi từ. Trạng thái thẻ
nhờ đó chuyển dịch \code{new} \( \to \) \code{learning} \( \to \)
\code{review} \( \to \) \code{mastered}, phù hợp với biểu đồ hoạt động đã
trình bày ở Chương~\ref{chuong:phan-tich-thiet-ke}.

\subsection{Khung tham chiếu trình độ ngôn ngữ CEFR}
\label{subsec:cefr-theory}

Hệ thống chuẩn hoá độ khó của từ vựng và trình độ người học theo
\textbf{Khung tham chiếu trình độ ngôn ngữ chung châu Âu} (Common European
Framework of Reference for Languages --- CEFR) do Hội đồng châu Âu ban hành
\cite{coe2001cefr}. CEFR chia năng lực ngôn ngữ thành sáu bậc từ
\( A1, A2 \) (sơ cấp), \( B1, B2 \) (trung cấp) đến \( C1, C2 \) (cao cấp).
Mỗi từ trong kho được gắn một bậc CEFR và mỗi người học khai báo trình độ
của mình; nhờ đó hệ thống lựa chọn từ và gợi ý bộ thẻ phù hợp với năng lực
hiện tại, đồng thời mô hình ngôn ngữ chấm điểm bậc CEFR mà một câu luyện
tập thể hiện.

\subsection{Kiến trúc REST và dịch vụ web}
\label{subsec:rest-theory}

Giao diện lập trình ứng dụng (API) của hệ thống tuân theo phong cách kiến
trúc \textbf{REST} (Representational State Transfer) do Fielding đề xuất
\cite{fielding2000rest}. Các ràng buộc cốt lõi của REST mà hệ thống áp dụng
gồm: mô hình \textit{client--server} tách bạch máy khách và máy chủ; tính
\textit{phi trạng thái} (stateless) --- mỗi yêu cầu tự mang đủ thông tin xác
thực (token), máy chủ không lưu trạng thái phiên; định danh tài nguyên bằng
\textit{URI} và thao tác qua các phương thức HTTP chuẩn (GET, POST, PATCH,
DELETE) với ngữ nghĩa thống nhất; và trao đổi dữ liệu dưới dạng
\textbf{JSON}. Tính phi trạng thái là điều kiện để mở rộng tầng API theo
chiều ngang (đáp ứng yêu cầu NFR-03 trong
Bảng~\ref{tab:nfr}). API còn được \textit{phiên bản hoá theo URI}
(tiền tố \code{/v1}) nhằm tiến hoá hợp đồng mà không phá vỡ máy khách cũ.

\subsection{Ánh xạ đối tượng--quan hệ (ORM)}
\label{subsec:orm-theory}

Để thu hẹp khoảng cách giữa mô hình hướng đối tượng trong mã nguồn và mô
hình quan hệ trong cơ sở dữ liệu, hệ thống dùng kỹ thuật \textbf{ánh xạ đối
tượng--quan hệ} (Object--Relational Mapping --- ORM). Mỗi bảng được khai
báo dưới dạng một lớp \textit{thực thể} (entity), mỗi dòng tương ứng một đối
tượng và các quan hệ khoá ngoại được biểu diễn bằng liên kết giữa các thực
thể. ORM giúp lập trình viên thao tác dữ liệu bằng mã có kiểu (typed) thay
vì viết câu lệnh SQL thủ công, đồng thời quản lý tiến hoá lược đồ bằng cơ
chế \textit{migration} có thể phiên bản hoá. Đánh đổi của ORM là một lớp
trừu tượng phụ và nguy cơ truy vấn kém hiệu quả nếu lạm dụng; hệ thống khắc
phục bằng cách đánh chỉ mục các truy vấn nóng và truy vấn tường minh khi cần.

\subsection{Mô hình ngôn ngữ lớn và sinh nội dung}
\label{subsec:llm-theory}

Các chức năng hỗ trợ bằng AI (làm giàu dữ liệu từ điển, chấm điểm câu luyện
tập) dựa trên \textbf{mô hình ngôn ngữ lớn} (Large Language Model --- LLM).
LLM là mạng nơ-ron kiến trúc \textit{Transformer} được huấn luyện trên khối
dữ liệu văn bản khổng lồ để dự đoán token kế tiếp, nhờ đó có khả năng sinh
văn bản mạch lạc, dịch và đánh giá ngôn ngữ theo chỉ dẫn (prompt). Hệ thống
sử dụng họ mô hình mở \textbf{Gemma} của Google DeepMind
\cite{gemma2025technical}. Trong nghiệp vụ, LLM được giao vai trò có kiểm
soát: sinh nghĩa, ví dụ, bản dịch theo định dạng yêu cầu và chấm điểm câu
theo \textit{rubric}; kết quả do AI sinh ra đối với kho hệ thống luôn phải
qua bước \textit{duyệt} của quản trị viên trước khi xuất bản, nhằm kiểm soát
chất lượng và hạn chế lỗi "ảo giác" (hallucination) của mô hình.

% =====================================================================
\section{Tổng quan công nghệ sử dụng}
\label{sec:tong-quan-cong-nghe}

Trên cơ sở lý thuyết ở mục~\ref{sec:co-so-ly-thuyet}, hệ thống lựa chọn ngăn
xếp công nghệ được tổ chức theo nhiều tầng. Bảng~\ref{tab:tech-overview}
tóm tắt công nghệ chủ chốt của từng tầng cùng vai trò của nó;
Hình~\ref{fig:tech-stack} minh hoạ ngăn xếp này một cách trực quan. Các mục
tiếp theo phân tích chi tiết từng nhóm công nghệ.

\begin{table}[h]
  \centering
  \small
  \renewcommand{\arraystretch}{1.3}
  \begin{tabular}{|p{3.0cm}|p{5.0cm}|p{6.4cm}|}
    \hline
    \textbf{Tầng} & \textbf{Công nghệ chính} & \textbf{Vai trò} \\
    \hline
    Frontend & Next.js, React, TypeScript, Tailwind CSS, shadcn/ui & Ứng dụng khách tiêu thụ API, dựng giao diện người dùng. \\
    \hline
    Backend & NestJS 11, TypeScript, TypeORM & Cung cấp API REST, chứa nghiệp vụ và truy cập dữ liệu. \\
    \hline
    Cơ sở dữ liệu & PostgreSQL & Lưu trữ dữ liệu quan hệ có ràng buộc toàn vẹn. \\
    \hline
    Hàng đợi \& nền & Redis, BullMQ & Hàng đợi công việc và worker xử lý bất đồng bộ. \\
    \hline
    Xác thực & JWT, Passport, bcrypt, OAuth/OIDC & Xác thực, phân quyền và bảo vệ tài khoản. \\
    \hline
    Trí tuệ nhân tạo & Gemma~3 (Google AI Studio), Edge TTS, vi dịch vụ chấm phát âm & Sinh nội dung, sinh âm thanh, chấm điểm luyện tập. \\
    \hline
    Hạ tầng phụ trợ & Cloudinary, SMTP, Docker & Lưu trữ media, gửi email, đóng gói triển khai. \\
    \hline
  \end{tabular}
  \caption{Tổng quan ngăn xếp công nghệ của hệ thống}
  \label{tab:tech-overview}
\end{table}

% --- PlantUML gợi ý cho Hình ngăn xếp công nghệ ----------------------
% @startuml
% rectangle "FRONTEND\nNext.js + React + TypeScript\nTailwind CSS + shadcn/ui" as FE
% rectangle "BACKEND (API)\nNestJS 11 + TypeScript\nController - Service - Repository (TypeORM)" as BE
% rectangle "XỬ LÝ NỀN\nWorker (BullMQ) + Redis" as WK
% database "PostgreSQL" as DB
% cloud "DỊCH VỤ AI / NGOÀI\nGemma 3, Edge TTS, Cloudinary,\nDV chấm phát âm, OAuth, SMTP" as EXT
% FE --> BE : HTTPS / JSON (REST)
% BE --> DB
% BE --> WK : enqueue
% WK --> DB
% WK --> EXT
% BE --> EXT
% @enduml
\figph{tech-stack.png}{0.8}{Ngăn xếp công nghệ của hệ thống theo tầng}{fig:tech-stack}

% =====================================================================
\section{Tổng quan kiến trúc hệ thống}
\label{sec:kien-truc-tong-quan}

Mục này điểm lại các \textit{kiểu kiến trúc} (architectural styles) làm nền
cho thiết kế đã trình bày chi tiết ở mục kiến trúc của
Chương~\ref{chuong:phan-tich-thiet-ke} (xem Hình~\ref{fig:kien-truc}).

\begin{itemize}
  \item \textbf{Monolith mô-đun hoá (modular monolith):} toàn bộ backend là
        một tiến trình triển khai duy nhất nhưng được chia thành các mô-đun
        nghiệp vụ tách bạch (xem Bảng~\ref{tab:modules}). So với kiến trúc
        vi dịch vụ (microservices), cách này đơn giản hơn khi triển khai và
        vận hành ở quy mô đồ án, đồng thời vẫn giữ ranh giới mô-đun rõ ràng
        để có thể tách dịch vụ về sau.
  \item \textbf{Kiến trúc phân tầng (layered):} trong mỗi mô-đun, mã nguồn
        chia theo các tầng \textit{Controller} (tiếp nhận HTTP, kiểm tra
        DTO) \( \to \) \textit{Service} (nghiệp vụ) \( \to \)
        \textit{Repository} (truy cập dữ liệu qua ORM). Việc tách tầng giúp
        kiểm thử và bảo trì độc lập từng mối quan tâm.
  \item \textbf{Client--server phi trạng thái (REST):} máy khách
        (Next.js) và máy chủ giao tiếp qua HTTP/JSON theo REST; máy chủ phi
        trạng thái nên có thể nhân bản để mở rộng theo chiều ngang.
  \item \textbf{Xử lý bất đồng bộ kiểu nhà sản xuất--người tiêu thụ
        (producer--consumer):} tiến trình API đóng vai trò nhà sản xuất, đẩy
        công việc nặng vào hàng đợi; tiến trình worker là người tiêu thụ,
        xử lý nền. Mẫu này giữ cho luồng yêu cầu chính luôn nhanh và ổn định.
\end{itemize}

Sự kết hợp bốn kiểu kiến trúc trên cho phép hệ thống vừa đơn giản khi triển
khai, vừa đáp ứng các yêu cầu phi chức năng về hiệu năng, khả năng mở rộng
và khả năng bảo trì trong Bảng~\ref{tab:nfr}.

% =====================================================================
\section{Công nghệ cơ sở dữ liệu}
\label{sec:cong-nghe-csdl}

Hệ thống dùng hệ quản trị cơ sở dữ liệu quan hệ \textbf{PostgreSQL}
\cite{postgresql2024}. PostgreSQL được chọn vì là mã nguồn mở, trưởng thành,
tuân thủ chuẩn SQL và hỗ trợ tốt các đặc tính mà thiết kế dữ liệu
(mục thiết kế CSDL của Chương~\ref{chuong:phan-tich-thiet-ke}) yêu cầu:

\begin{itemize}
  \item \textbf{Giao dịch ACID:} bảo đảm tính nguyên tử cho các thao tác
        ghi nhiều bảng (ví dụ tạo phiên học, cập nhật tiến độ và ghi nhật
        ký), đáp ứng yêu cầu độ tin cậy.
  \item \textbf{Khoá chính \code{uuid}:} dùng kiểu định danh duy nhất toàn
        cục cho mọi bảng, thuận lợi cho phân tán và tránh lộ số lượng bản ghi.
  \item \textbf{Kiểu \code{timestamptz}:} lưu mốc thời gian kèm múi giờ,
        bảo đảm nhất quán cho thống kê chuỗi ngày học và lịch ôn tập.
  \item \textbf{Ràng buộc toàn vẹn:} khoá ngoại, ràng buộc duy nhất và quy
        tắc xoá lan truyền (\code{CASCADE} / \code{SET NULL}) giữ dữ liệu
        nhất quán giữa từ vựng, bộ thẻ và tiến độ.
  \item \textbf{Chỉ mục:} đánh chỉ mục các cột truy vấn nóng (lemma, trạng
        thái ôn tập, thời điểm đến hạn) để bảo đảm hiệu năng khi dữ liệu lớn.
\end{itemize}

Bên cạnh PostgreSQL, hệ thống dùng \textbf{Redis} làm kho khoá--giá trị
trong bộ nhớ, đóng vai trò \textit{nơi lưu trữ hàng đợi} cho BullMQ
\cite{redis2024}. Việc truy cập cơ sở dữ liệu từ phía backend được thực hiện
qua \textbf{TypeORM} --- một thư viện ORM cho TypeScript/Node.js
\cite{typeorm2024} --- sẽ được trình bày ở mục~\ref{sec:cong-nghe-backend}.
Lược đồ cơ sở dữ liệu được quản lý bằng các \textit{migration} đặt trong thư
mục mã nguồn và chạy qua công cụ dòng lệnh của TypeORM, giúp tiến hoá lược
đồ một cách có kiểm soát và tái lập được giữa các môi trường.

% =====================================================================
\section{Công nghệ Frontend}
\label{sec:cong-nghe-frontend}

Ứng dụng khách (frontend) là một dự án độc lập, đóng vai trò bên tiêu thụ
API REST của backend. Frontend được xây dựng bằng \textbf{Next.js 16} theo
mô hình \textit{App Router} --- framework React hỗ trợ kết xuất phía máy chủ
(Server-Side Rendering --- SSR), \textit{thành phần phía máy chủ} (Server
Components) và \textit{hành động phía máy chủ} (Server Actions)
\cite{nextjs2024} --- trên nền thư viện giao diện \textbf{React~19}
\cite{react2024} và ngôn ngữ \textbf{TypeScript} (chế độ \textit{strict}).
Việc dùng chung TypeScript ở cả frontend và backend giúp thống nhất kiểu dữ
liệu của hợp đồng API giữa hai phía. Phần tạo kiểu giao diện dùng
\textbf{Tailwind CSS v4} --- framework CSS theo hướng tiện ích (utility-first),
cấu hình bằng CSS \cite{tailwind2024} --- kết hợp bộ thành phần
\textbf{shadcn/ui} dựng trên các primitive khả truy cập của \textbf{Base UI}.
Ràng buộc và kiểm tra dữ liệu biểu mẫu dùng \code{zod} \cite{zod2024}, được
dùng lại giữa kiểm tra phía máy khách và Server Action. Bảng~\ref{tab:frontend-libs}
liệt kê các thư viện frontend chính.

\begin{table}[h]
  \centering
  \small
  \renewcommand{\arraystretch}{1.3}
  \begin{tabular}{|p{4.2cm}|p{10.4cm}|}
    \hline
    \textbf{Công nghệ / thư viện} & \textbf{Vai trò} \\
    \hline
    Next.js 16 (App Router), React 19, TypeScript & Khung ứng dụng theo App Router; Server Components/Server Actions, định tuyến và kết xuất trang. \\
    \hline
    Tailwind CSS v4, shadcn/ui, Base UI & Hệ thống tạo kiểu tiện ích và bộ thành phần giao diện khả truy cập, tái sử dụng. \\
    \hline
    \code{zod} & Khai báo và kiểm tra ràng buộc dữ liệu (biểu mẫu, phản hồi API) phía máy khách. \\
    \hline
    \code{lucide-react}, \code{sonner} & Bộ biểu tượng và thông báo nổi (toast) cho phản hồi thao tác. \\
    \hline
    \code{next-themes}, \code{input-otp} & Hỗ trợ chế độ sáng/tối và ô nhập mã OTP (xác minh email). \\
    \hline
  \end{tabular}
  \caption{Các công nghệ và thư viện frontend chính}
  \label{tab:frontend-libs}
\end{table}

Mục này chỉ giới thiệu \textit{ngăn xếp công nghệ} frontend; quá trình hiện
thực hoá tầng giao diện (cấu trúc App Router, tầng truy cập API phía máy chủ,
quản lý phiên và động cơ phiên học phía máy khách) được trình bày ở
mục~\ref{sec:hien-thuc-frontend}, Chương~\ref{chuong:hien-thuc-hoa}. Trong mọi
trường hợp, frontend giao tiếp với backend hoàn toàn qua các API đã đặc tả,
không truy cập trực tiếp cơ sở dữ liệu.

% =====================================================================
\section{Công nghệ Backend}
\label{sec:cong-nghe-backend}

Backend được xây dựng bằng \textbf{NestJS}, một framework Node.js cho
TypeScript theo hướng mô-đun, dựa trên các nguyên lý \textit{tiêm phụ thuộc}
(dependency injection) và lập trình hướng khía cạnh \cite{nestjs2024}.
NestJS phù hợp với kiến trúc phân tầng đã nêu ở
mục~\ref{sec:kien-truc-tong-quan}: mỗi mô-đun đóng gói controller, service
và provider riêng; các \textit{guard}, \textit{pipe} và
\textit{interceptor} xử lý những mối quan tâm xuyên suốt (xác thực, kiểm tra
đầu vào, biến đổi phản hồi) một cách khai báo. Toàn bộ mã nguồn viết bằng
\textbf{TypeScript} ở chế độ \textit{strict}, giúp phát hiện lỗi kiểu ngay
khi biên dịch. Bảng~\ref{tab:backend-libs} liệt kê các thành phần backend
chính.

\begin{table}[h]
  \centering
  \small
  \renewcommand{\arraystretch}{1.3}
  \begin{tabular}{|p{4.6cm}|p{10.0cm}|}
    \hline
    \textbf{Công nghệ / thư viện} & \textbf{Vai trò} \\
    \hline
    NestJS 11 & Framework backend: tổ chức mô-đun, định tuyến, tiêm phụ thuộc, phiên bản hoá API theo URI. \\
    \hline
    TypeORM & Ánh xạ đối tượng--quan hệ tới PostgreSQL; quản lý thực thể và migration. \\
    \hline
    \code{class-validator}, \code{class-transformer} & Kiểm tra ràng buộc DTO và biến đổi dữ liệu vào/ra qua \code{ValidationPipe} toàn cục (loại bỏ trường lạ). \\
    \hline
    BullMQ + Redis & Hàng đợi công việc và worker xử lý nền (làm giàu từ vựng, sinh âm thanh, chấm điểm luyện tập). \\
    \hline
    \code{@nestjs/throttler} & Giới hạn tần suất (rate limiting) cho các điểm cuối nhạy cảm như đăng nhập. \\
    \hline
    \code{nodemailer} & Gửi email (mã xác minh) qua SMTP. \\
    \hline
    \code{pdf-parse}, \code{xlsx} & Trích xuất văn bản từ tệp PDF và bảng tính phục vụ nhập từ vựng hàng loạt. \\
    \hline
  \end{tabular}
  \caption{Các công nghệ và thư viện backend chính}
  \label{tab:backend-libs}
\end{table}

Phần xử lý bất đồng bộ được hiện thực bằng \textbf{BullMQ} --- thư viện hàng
đợi công việc chạy trên Redis \cite{bullmq2024}. Tiến trình API và tiến
trình worker chia sẻ chung cơ sở dữ liệu và hàng đợi nhưng được khởi chạy
độc lập (lệnh \code{start} và \code{start:worker}), phản ánh đúng mô hình
producer--consumer ở mục~\ref{sec:kien-truc-tong-quan}.

% =====================================================================
\section{Cơ chế xác thực và phân quyền}
\label{sec:xac-thuc-phan-quyen}

Phân hệ tài khoản áp dụng nhiều lớp kỹ thuật bảo mật, đáp ứng yêu cầu
NFR-01 trong Bảng~\ref{tab:nfr}.

\subsection{Băm mật khẩu}
\label{subsec:bcrypt}

Mật khẩu người dùng không bao giờ lưu ở dạng văn bản thuần mà được băm bằng
\textbf{bcrypt} --- hàm băm mật khẩu thích nghi, có \textit{muối} (salt) và
hệ số chi phí điều chỉnh được để chống tấn công vét cạn
\cite{provos1999bcrypt}. Khi đăng nhập, hệ thống so khớp mật khẩu nhập vào
với bản băm đã lưu thay vì giải mã.

\subsection{Xác thực bằng JSON Web Token}
\label{subsec:jwt}

Sau khi xác thực thành công, hệ thống cấp một cặp token theo chuẩn
\textbf{JSON Web Token (JWT)} \cite{rfc7519}: một \textit{access token} sống
ngắn để gọi API và một \textit{refresh token} sống dài để làm mới phiên.
JWT là chuỗi tự chứa, được ký số nên máy chủ có thể xác minh tính toàn vẹn
mà không cần tra cứu trạng thái phiên --- đúng tinh thần phi trạng thái của
REST (mục~\ref{subsec:rest-theory}). Thư viện \textbf{Passport} (qua
\code{passport-jwt}) đảm nhiệm việc trích và xác minh token trong mỗi yêu
cầu. Để có thể \textit{thu hồi}, refresh token được lưu dưới dạng băm trong
bảng \code{refresh_tokens} và bị vô hiệu khi người dùng đăng xuất.

\subsection{Phân quyền theo vai trò và quyền sở hữu}
\label{subsec:rbac}

Việc cấp quyền dựa trên hai cơ chế bổ trợ:

\begin{itemize}
  \item \textbf{Theo vai trò (role-based):} mỗi tài khoản mang vai trò
        \code{user} hoặc \code{admin}; một \textit{guard} kiểm tra vai trò
        để chặn các thao tác quản trị đối với người dùng thường.
  \item \textbf{Theo quyền sở hữu (ownership-based):} với tài nguyên cá
        nhân (từ vựng riêng, bộ thẻ, tiến độ), hệ thống kiểm tra người gọi
        có đúng là chủ sở hữu hay không trước khi cho phép truy cập.
\end{itemize}

\subsection{Đăng nhập qua nhà cung cấp bên thứ ba (OAuth/OIDC)}
\label{subsec:oauth}

Ngoài email/mật khẩu, hệ thống cho phép đăng nhập qua \textbf{Google},
\textbf{Apple} và \textbf{GitHub} bằng giao thức \textbf{OAuth~2.0} và lớp
nhận dạng \textbf{OpenID Connect} (OIDC) \cite{rfc6749}. Token danh tính do
nhà cung cấp cấp được hệ thống xác minh, sau đó liên kết với một bản ghi
trong bảng \code{user_identities}; nhờ vậy một người dùng có thể đăng nhập
bằng nhiều phương thức mà vẫn chung một tài khoản.

\subsection{Chấm điểm phiên học phi trạng thái bằng HMAC}
\label{subsec:hmac}

Một điểm đặc thù về bảo mật là cơ chế \textbf{chấm điểm phi trạng thái} cho
phiên học. Thay vì lưu trạng thái phiên trên máy chủ, mỗi câu hỏi khi sinh
ra được \textit{ký} bằng \textbf{HMAC} (Hash-based Message Authentication
Code) \cite{rfc2104} dựa trên một khoá bí mật của máy chủ. Khi người học nộp
đáp án, máy chủ xác minh lại chữ ký để bảo đảm câu hỏi không bị giả mạo hay
hết hạn, rồi mới suy lại đáp án đúng và chấm điểm. Cách làm này vừa giữ tính
phi trạng thái (thuận lợi cho mở rộng ngang) vừa chống gian lận.

% =====================================================================
\section{Công nghệ trí tuệ nhân tạo}
\label{sec:cong-nghe-ai}

Các chức năng thông minh của hệ thống dựa trên một số dịch vụ AI được tích
hợp qua tầng worker xử lý nền. Bảng~\ref{tab:ai-services} tổng hợp các dịch
vụ này.

\begin{table}[h]
  \centering
  \small
  \renewcommand{\arraystretch}{1.3}
  \begin{tabular}{|p{3.6cm}|p{4.0cm}|p{6.8cm}|}
    \hline
    \textbf{Dịch vụ} & \textbf{Loại} & \textbf{Vai trò trong hệ thống} \\
    \hline
    Gemma~3 (Google AI Studio) & Mô hình ngôn ngữ lớn (LLM) & Sinh nghĩa, ví dụ, bản dịch; chấm điểm câu luyện viết theo rubric kèm bậc CEFR. \\
    \hline
    Microsoft Edge TTS & Tổng hợp giọng nói (TTS) & Sinh âm thanh phát âm tự động cho từ vựng. \\
    \hline
    Vi dịch vụ chấm phát âm & Mô hình chấm điểm âm vị & Nhận bản ghi âm, trả điểm tổng và điểm từng phoneme. \\
    \hline
  \end{tabular}
  \caption{Các dịch vụ trí tuệ nhân tạo được tích hợp}
  \label{tab:ai-services}
\end{table}

\subsection{Sinh nội dung và chấm điểm bằng Gemma~3}
\label{subsec:gemma}

Hệ thống gọi mô hình \textbf{Gemma~3} qua nền tảng Google AI Studio
\cite{gemma2025technical}. Mô hình được sử dụng cho hai nhóm tác vụ:
(i)~\textit{làm giàu dữ liệu từ điển} --- từ một lemma, sinh các nghĩa, câu
ví dụ và bản dịch theo định dạng có cấu trúc; và (ii)~\textit{chấm điểm câu
luyện tập} --- đánh giá câu người học đặt theo thang điểm 0--100 kèm bậc
CEFR và nhận xét. Mọi lời gọi LLM đều là tác vụ nặng nên được đẩy sang
worker xử lý bất đồng bộ và bị giới hạn theo hạn mức ngày để bảo vệ khoá API
dùng chung.

\subsection{Tổng hợp giọng nói và chấm điểm phát âm}
\label{subsec:tts-pron}

Âm thanh phát âm của từ vựng được sinh tự động bằng \textbf{Microsoft Edge
TTS} (qua thư viện \code{msedge-tts}) rồi tải lên kho media. Chức năng luyện
phát âm được phục vụ bởi một \textbf{vi dịch vụ} riêng viết bằng Python
(FastAPI): backend đóng vai trò proxy mỏng, chuyển tiếp bản ghi âm của người
học tới vi dịch vụ để nhận điểm tổng và điểm từng âm vị, sau đó lưu lại lịch
sử luyện tập. Việc tách phần chấm phát âm thành dịch vụ độc lập giúp cô lập
các phụ thuộc nặng về xử lý tín hiệu/học máy khỏi tiến trình backend chính.

% =====================================================================
\section{Công cụ phát triển}
\label{sec:cong-cu-phat-trien}

Quá trình phát triển sử dụng bộ công cụ trong Bảng~\ref{tab:dev-tools} nhằm
bảo đảm chất lượng mã nguồn, khả năng tái lập môi trường và năng suất làm
việc.

\begin{table}[h]
  \centering
  \small
  \renewcommand{\arraystretch}{1.3}
  \begin{tabular}{|p{4.0cm}|p{10.6cm}|}
    \hline
    \textbf{Công cụ} & \textbf{Mục đích} \\
    \hline
    Git \& GitHub & Quản lý phiên bản mã nguồn và cộng tác qua pull request. \\
    \hline
    Docker \& Docker Compose & Đóng gói và khởi chạy nhất quán PostgreSQL, Redis, pgAdmin phục vụ phát triển và triển khai \cite{docker2024}. \\
    \hline
    NestJS CLI & Sinh khung mô-đun/controller/service và biên dịch, chạy ứng dụng. \\
    \hline
    TypeORM CLI & Sinh, chạy và hoàn tác migration của lược đồ cơ sở dữ liệu. \\
    \hline
    ESLint \& Prettier & Phân tích tĩnh và định dạng mã TypeScript theo quy ước thống nhất. \\
    \hline
    Jest \& Supertest & Kiểm thử đơn vị và kiểm thử tích hợp (end-to-end) cho API. \\
    \hline
    pgAdmin & Quản trị và truy vấn trực quan cơ sở dữ liệu PostgreSQL. \\
    \hline
    Overleaf (\LaTeX) & Biên soạn báo cáo đồ án theo mẫu của nhà trường. \\
    \hline
  \end{tabular}
  \caption{Các công cụ phát triển sử dụng trong đề tài}
  \label{tab:dev-tools}
\end{table}

Mã nguồn được quản lý bằng \textbf{Git} và lưu trên GitHub; mỗi tính năng
hoặc sửa lỗi được phát triển trên một nhánh riêng rồi hợp nhất qua pull
request. Môi trường phụ thuộc (PostgreSQL, Redis, pgAdmin) được mô tả bằng
\textbf{Docker Compose} \cite{docker2024}, giúp tái lập nhanh và đồng nhất
giữa các máy. Chất lượng mã được kiểm soát bằng kiểm tra kiểu của TypeScript,
\textbf{ESLint} và \textbf{Prettier}; phần kiểm thử dùng \textbf{Jest} cho
kiểm thử đơn vị và \textbf{Supertest} cho kiểm thử tích hợp các điểm cuối
API. Cuối cùng, báo cáo đồ án được biên soạn bằng \textbf{\LaTeX} trên
Overleaf theo đúng mẫu quy định.

% =====================================================================
\section{Kết chương}
\label{sec:ket-chuong-3}

Chương này đã trình bày các cơ sở lý thuyết --- ôn tập ngắt quãng và thuật
toán SM-2, khung trình độ CEFR, kiến trúc REST, ánh xạ đối tượng--quan hệ và
mô hình ngôn ngữ lớn --- cùng ngăn xếp công nghệ được lựa chọn cho từng tầng
của hệ thống, từ cơ sở dữ liệu, frontend, backend, cơ chế xác thực và phân
quyền, các dịch vụ trí tuệ nhân tạo đến bộ công cụ phát triển. Những lựa
chọn này trực tiếp phục vụ các quyết định thiết kế ở
Chương~\ref{chuong:phan-tich-thiet-ke} và là tiền đề cho quá trình hiện
thực hoá và kiểm thử được trình bày ở chương tiếp theo.

%  vào main.tex (sau Chương 2).
%
%  GHI CHÚ VỀ HÌNH VẼ: dùng macro \figph{<tên-file>}{<tỉ-lệ-rộng>}{<chú-thích>}{<nhãn>}
%  (đã khai báo ở Chương 2, khai báo lại bằng \providecommand để an toàn nếu
%  biên dịch độc lập). Macro tự chèn ảnh thật nếu tệp tồn tại trong images/,
%  ngược lại hiển thị khung giữ chỗ để báo cáo vẫn biên dịch được.
%
%  KHÔNG hardcode số chương/mục/bảng/hình/công thức — luôn dùng \label + \ref / \cite.
% =====================================================================

% --- Macro tiện ích (idempotent — không định nghĩa lại nếu đã có) ------
\providecommand{\code}[1]{\texttt{\detokenize{#1}}} % in định danh snake_case an toàn
\providecommand{\figph}[4]{%
  \begin{figure}[htbp]
    \centering
    \IfFileExists{images/#1}%
      {\includegraphics[width=#2\textwidth,height=0.82\textheight,keepaspectratio]{#1}}%
      {\setlength{\fboxsep}{10pt}\fbox{\begin{minipage}[c][3.6cm][c]{#2\textwidth}\centering\itshape Sơ đồ minh hoạ --- chèn tệp ảnh \code{images/#1} vào thư mục \code{images/}.\end{minipage}}}%
    \caption{#3}%
    \label{#4}%
  \end{figure}}

\chapter{CƠ SỞ LÝ THUYẾT VÀ CÔNG NGHỆ SỬ DỤNG}
\label{chuong:co-so-ly-thuyet}

Chương~\ref{chuong:phan-tich-thiet-ke} đã trình bày kết quả phân tích và
thiết kế hệ thống. Chương này hệ thống hoá những \textit{cơ sở lý thuyết}
và \textit{công nghệ} làm nền tảng cho các quyết định thiết kế và hiện thực
hoá đó, qua đó lý giải vì sao mỗi công nghệ được lựa chọn. Trước hết,
mục~\ref{sec:co-so-ly-thuyet} trình bày các cơ sở lý thuyết cốt lõi: ôn tập
ngắt quãng và thuật toán SM-2, khung trình độ CEFR, kiến trúc REST, ánh xạ
đối tượng--quan hệ và mô hình ngôn ngữ lớn. Tiếp theo,
mục~\ref{sec:tong-quan-cong-nghe} tổng quan ngăn xếp công nghệ và
mục~\ref{sec:kien-truc-tong-quan} phân tích các kiểu kiến trúc được áp dụng.
Các mục~\ref{sec:cong-nghe-csdl}--\ref{sec:cong-cu-phat-trien} đi sâu lần
lượt vào công nghệ cơ sở dữ liệu, frontend, backend, cơ chế xác thực và
phân quyền, công nghệ trí tuệ nhân tạo và bộ công cụ phát triển.

% =====================================================================
\section{Cơ sở lý thuyết}
\label{sec:co-so-ly-thuyet}

\subsection{Ôn tập ngắt quãng và đường cong lãng quên}
\label{subsec:srs-theory}

Cơ sở khoa học của động cơ học trong hệ thống là hiện tượng
\textit{đường cong lãng quên} (forgetting curve) do Ebbinghaus phát hiện:
khả năng ghi nhớ một thông tin mới suy giảm theo hàm mũ theo thời gian nếu
không được củng cố \cite{ebbinghaus1913memory}. Một cách mô hình hoá đơn
giản, mức ghi nhớ \( R \) còn lại sau khoảng thời gian \( t \) kể từ lần
học gần nhất được biểu diễn qua công thức~\eqref{eq:forgetting}:

\begin{equation}
  \label{eq:forgetting}
  R(t) = e^{-\,t/S},
\end{equation}

trong đó \( S \) là \textit{độ bền trí nhớ} (memory stability) --- càng ôn
tập đúng lúc, \( S \) càng lớn và thông tin được giữ lâu hơn. Hệ quả thực
tiễn là: thay vì học nhồi nhét, người học nên ôn lại mỗi từ tại những
\textit{thời điểm đến hạn} được giãn dần. Các nghiên cứu về hiệu ứng giãn
cách (spacing effect) khẳng định ôn tập ngắt quãng cải thiện rõ rệt khả
năng ghi nhớ dài hạn so với học dồn \cite{cepeda2006distributed}. Đây là
nền tảng lý thuyết cho \textit{hệ thống ôn tập ngắt quãng} (Spaced
Repetition System --- SRS) mà hệ thống hiện thực hoá.

\subsection{Thuật toán lập lịch ôn tập SM-2}
\label{subsec:sm2-theory}

Để quyết định \textit{khi nào} ôn lại mỗi từ, hệ thống áp dụng thuật toán
\textbf{SM-2} có nguồn gốc từ dự án SuperMemo \cite{wozniak1990optimization}.
Mỗi cặp (người dùng, từ) được mô tả bằng ba đại lượng: hệ số dễ nhớ
\textit{easiness factor} \( EF \) (khởi tạo \( 2{,}5 \)), số lần ôn đúng
liên tiếp \( n \) và khoảng cách ôn tập \( I \) (tính theo ngày). Sau mỗi
lượt ôn, người học (ở đây là kết quả trả lời trong phiên) được quy về một
\textit{điểm chất lượng} \( q \in \{0,1,\dots,5\} \).

Hệ số \( EF \) được cập nhật theo công thức~\eqref{eq:sm2-ef}, có chặn dưới
\( 1{,}3 \) để khoảng ôn tập không co lại quá nhanh:

\begin{equation}
  \label{eq:sm2-ef}
  EF' = EF + \big(0{,}1 - (5-q)\,(0{,}08 + (5-q)\cdot 0{,}02)\big),
  \qquad EF' \geq 1{,}3 .
\end{equation}

Khoảng cách ôn tập kế tiếp được tính theo công thức~\eqref{eq:sm2-interval}:

\begin{equation}
  \label{eq:sm2-interval}
  I(n) =
  \begin{cases}
    1, & n = 1,\\[2pt]
    6, & n = 2,\\[2pt]
    \big\lceil I(n-1)\times EF' \big\rceil, & n > 2 .
  \end{cases}
\end{equation}

Nếu người học trả lời sai (\( q < 3 \)), bộ đếm \( n \) được đặt lại về 0 và
từ phải học lại từ đầu; ngược lại \( n \) tăng và khoảng ôn tập giãn dần
theo \( EF' \). Trong hệ thống, SM-2 được \textbf{mở rộng bằng các bước học}
(learning steps): một từ mới phải vượt qua một chuỗi câu hỏi từ dễ đến khó
trong phiên trước khi được giao cho SM-2 lập lịch dài hạn; điểm chất lượng
\( q \) được suy ra từ độ chính xác ở bước cuối của mỗi từ. Trạng thái thẻ
nhờ đó chuyển dịch \code{new} \( \to \) \code{learning} \( \to \)
\code{review} \( \to \) \code{mastered}, phù hợp với biểu đồ hoạt động đã
trình bày ở Chương~\ref{chuong:phan-tich-thiet-ke}.

\subsection{Khung tham chiếu trình độ ngôn ngữ CEFR}
\label{subsec:cefr-theory}

Hệ thống chuẩn hoá độ khó của từ vựng và trình độ người học theo
\textbf{Khung tham chiếu trình độ ngôn ngữ chung châu Âu} (Common European
Framework of Reference for Languages --- CEFR) do Hội đồng châu Âu ban hành
\cite{coe2001cefr}. CEFR chia năng lực ngôn ngữ thành sáu bậc từ
\( A1, A2 \) (sơ cấp), \( B1, B2 \) (trung cấp) đến \( C1, C2 \) (cao cấp).
Mỗi từ trong kho được gắn một bậc CEFR và mỗi người học khai báo trình độ
của mình; nhờ đó hệ thống lựa chọn từ và gợi ý bộ thẻ phù hợp với năng lực
hiện tại, đồng thời mô hình ngôn ngữ chấm điểm bậc CEFR mà một câu luyện
tập thể hiện.

\subsection{Kiến trúc REST và dịch vụ web}
\label{subsec:rest-theory}

Giao diện lập trình ứng dụng (API) của hệ thống tuân theo phong cách kiến
trúc \textbf{REST} (Representational State Transfer) do Fielding đề xuất
\cite{fielding2000rest}. Các ràng buộc cốt lõi của REST mà hệ thống áp dụng
gồm: mô hình \textit{client--server} tách bạch máy khách và máy chủ; tính
\textit{phi trạng thái} (stateless) --- mỗi yêu cầu tự mang đủ thông tin xác
thực (token), máy chủ không lưu trạng thái phiên; định danh tài nguyên bằng
\textit{URI} và thao tác qua các phương thức HTTP chuẩn (GET, POST, PATCH,
DELETE) với ngữ nghĩa thống nhất; và trao đổi dữ liệu dưới dạng
\textbf{JSON}. Tính phi trạng thái là điều kiện để mở rộng tầng API theo
chiều ngang (đáp ứng yêu cầu NFR-03 trong
Bảng~\ref{tab:nfr}). API còn được \textit{phiên bản hoá theo URI}
(tiền tố \code{/v1}) nhằm tiến hoá hợp đồng mà không phá vỡ máy khách cũ.

\subsection{Ánh xạ đối tượng--quan hệ (ORM)}
\label{subsec:orm-theory}

Để thu hẹp khoảng cách giữa mô hình hướng đối tượng trong mã nguồn và mô
hình quan hệ trong cơ sở dữ liệu, hệ thống dùng kỹ thuật \textbf{ánh xạ đối
tượng--quan hệ} (Object--Relational Mapping --- ORM). Mỗi bảng được khai
báo dưới dạng một lớp \textit{thực thể} (entity), mỗi dòng tương ứng một đối
tượng và các quan hệ khoá ngoại được biểu diễn bằng liên kết giữa các thực
thể. ORM giúp lập trình viên thao tác dữ liệu bằng mã có kiểu (typed) thay
vì viết câu lệnh SQL thủ công, đồng thời quản lý tiến hoá lược đồ bằng cơ
chế \textit{migration} có thể phiên bản hoá. Đánh đổi của ORM là một lớp
trừu tượng phụ và nguy cơ truy vấn kém hiệu quả nếu lạm dụng; hệ thống khắc
phục bằng cách đánh chỉ mục các truy vấn nóng và truy vấn tường minh khi cần.

\subsection{Mô hình ngôn ngữ lớn và sinh nội dung}
\label{subsec:llm-theory}

Các chức năng hỗ trợ bằng AI (làm giàu dữ liệu từ điển, chấm điểm câu luyện
tập) dựa trên \textbf{mô hình ngôn ngữ lớn} (Large Language Model --- LLM).
LLM là mạng nơ-ron kiến trúc \textit{Transformer} được huấn luyện trên khối
dữ liệu văn bản khổng lồ để dự đoán token kế tiếp, nhờ đó có khả năng sinh
văn bản mạch lạc, dịch và đánh giá ngôn ngữ theo chỉ dẫn (prompt). Hệ thống
sử dụng họ mô hình mở \textbf{Gemma} của Google DeepMind
\cite{gemma2025technical}. Trong nghiệp vụ, LLM được giao vai trò có kiểm
soát: sinh nghĩa, ví dụ, bản dịch theo định dạng yêu cầu và chấm điểm câu
theo \textit{rubric}; kết quả do AI sinh ra đối với kho hệ thống luôn phải
qua bước \textit{duyệt} của quản trị viên trước khi xuất bản, nhằm kiểm soát
chất lượng và hạn chế lỗi "ảo giác" (hallucination) của mô hình.

% =====================================================================
\section{Tổng quan công nghệ sử dụng}
\label{sec:tong-quan-cong-nghe}

Trên cơ sở lý thuyết ở mục~\ref{sec:co-so-ly-thuyet}, hệ thống lựa chọn ngăn
xếp công nghệ được tổ chức theo nhiều tầng. Bảng~\ref{tab:tech-overview}
tóm tắt công nghệ chủ chốt của từng tầng cùng vai trò của nó;
Hình~\ref{fig:tech-stack} minh hoạ ngăn xếp này một cách trực quan. Các mục
tiếp theo phân tích chi tiết từng nhóm công nghệ.

\begin{table}[h]
  \centering
  \small
  \renewcommand{\arraystretch}{1.3}
  \begin{tabular}{|p{3.0cm}|p{5.0cm}|p{6.4cm}|}
    \hline
    \textbf{Tầng} & \textbf{Công nghệ chính} & \textbf{Vai trò} \\
    \hline
    Frontend & Next.js, React, TypeScript, Tailwind CSS, shadcn/ui & Ứng dụng khách tiêu thụ API, dựng giao diện người dùng. \\
    \hline
    Backend & NestJS 11, TypeScript, TypeORM & Cung cấp API REST, chứa nghiệp vụ và truy cập dữ liệu. \\
    \hline
    Cơ sở dữ liệu & PostgreSQL & Lưu trữ dữ liệu quan hệ có ràng buộc toàn vẹn. \\
    \hline
    Hàng đợi \& nền & Redis, BullMQ & Hàng đợi công việc và worker xử lý bất đồng bộ. \\
    \hline
    Xác thực & JWT, Passport, bcrypt, OAuth/OIDC & Xác thực, phân quyền và bảo vệ tài khoản. \\
    \hline
    Trí tuệ nhân tạo & Gemma~3 (Google AI Studio), Edge TTS, vi dịch vụ chấm phát âm & Sinh nội dung, sinh âm thanh, chấm điểm luyện tập. \\
    \hline
    Hạ tầng phụ trợ & Cloudinary, SMTP, Docker & Lưu trữ media, gửi email, đóng gói triển khai. \\
    \hline
  \end{tabular}
  \caption{Tổng quan ngăn xếp công nghệ của hệ thống}
  \label{tab:tech-overview}
\end{table}

% --- PlantUML gợi ý cho Hình ngăn xếp công nghệ ----------------------
% @startuml
% rectangle "FRONTEND\nNext.js + React + TypeScript\nTailwind CSS + shadcn/ui" as FE
% rectangle "BACKEND (API)\nNestJS 11 + TypeScript\nController - Service - Repository (TypeORM)" as BE
% rectangle "XỬ LÝ NỀN\nWorker (BullMQ) + Redis" as WK
% database "PostgreSQL" as DB
% cloud "DỊCH VỤ AI / NGOÀI\nGemma 3, Edge TTS, Cloudinary,\nDV chấm phát âm, OAuth, SMTP" as EXT
% FE --> BE : HTTPS / JSON (REST)
% BE --> DB
% BE --> WK : enqueue
% WK --> DB
% WK --> EXT
% BE --> EXT
% @enduml
\figph{tech-stack.png}{0.8}{Ngăn xếp công nghệ của hệ thống theo tầng}{fig:tech-stack}

% =====================================================================
\section{Tổng quan kiến trúc hệ thống}
\label{sec:kien-truc-tong-quan}

Mục này điểm lại các \textit{kiểu kiến trúc} (architectural styles) làm nền
cho thiết kế đã trình bày chi tiết ở mục kiến trúc của
Chương~\ref{chuong:phan-tich-thiet-ke} (xem Hình~\ref{fig:kien-truc}).

\begin{itemize}
  \item \textbf{Monolith mô-đun hoá (modular monolith):} toàn bộ backend là
        một tiến trình triển khai duy nhất nhưng được chia thành các mô-đun
        nghiệp vụ tách bạch (xem Bảng~\ref{tab:modules}). So với kiến trúc
        vi dịch vụ (microservices), cách này đơn giản hơn khi triển khai và
        vận hành ở quy mô đồ án, đồng thời vẫn giữ ranh giới mô-đun rõ ràng
        để có thể tách dịch vụ về sau.
  \item \textbf{Kiến trúc phân tầng (layered):} trong mỗi mô-đun, mã nguồn
        chia theo các tầng \textit{Controller} (tiếp nhận HTTP, kiểm tra
        DTO) \( \to \) \textit{Service} (nghiệp vụ) \( \to \)
        \textit{Repository} (truy cập dữ liệu qua ORM). Việc tách tầng giúp
        kiểm thử và bảo trì độc lập từng mối quan tâm.
  \item \textbf{Client--server phi trạng thái (REST):} máy khách
        (Next.js) và máy chủ giao tiếp qua HTTP/JSON theo REST; máy chủ phi
        trạng thái nên có thể nhân bản để mở rộng theo chiều ngang.
  \item \textbf{Xử lý bất đồng bộ kiểu nhà sản xuất--người tiêu thụ
        (producer--consumer):} tiến trình API đóng vai trò nhà sản xuất, đẩy
        công việc nặng vào hàng đợi; tiến trình worker là người tiêu thụ,
        xử lý nền. Mẫu này giữ cho luồng yêu cầu chính luôn nhanh và ổn định.
\end{itemize}

Sự kết hợp bốn kiểu kiến trúc trên cho phép hệ thống vừa đơn giản khi triển
khai, vừa đáp ứng các yêu cầu phi chức năng về hiệu năng, khả năng mở rộng
và khả năng bảo trì trong Bảng~\ref{tab:nfr}.

% =====================================================================
\section{Công nghệ cơ sở dữ liệu}
\label{sec:cong-nghe-csdl}

Hệ thống dùng hệ quản trị cơ sở dữ liệu quan hệ \textbf{PostgreSQL}
\cite{postgresql2024}. PostgreSQL được chọn vì là mã nguồn mở, trưởng thành,
tuân thủ chuẩn SQL và hỗ trợ tốt các đặc tính mà thiết kế dữ liệu
(mục thiết kế CSDL của Chương~\ref{chuong:phan-tich-thiet-ke}) yêu cầu:

\begin{itemize}
  \item \textbf{Giao dịch ACID:} bảo đảm tính nguyên tử cho các thao tác
        ghi nhiều bảng (ví dụ tạo phiên học, cập nhật tiến độ và ghi nhật
        ký), đáp ứng yêu cầu độ tin cậy.
  \item \textbf{Khoá chính \code{uuid}:} dùng kiểu định danh duy nhất toàn
        cục cho mọi bảng, thuận lợi cho phân tán và tránh lộ số lượng bản ghi.
  \item \textbf{Kiểu \code{timestamptz}:} lưu mốc thời gian kèm múi giờ,
        bảo đảm nhất quán cho thống kê chuỗi ngày học và lịch ôn tập.
  \item \textbf{Ràng buộc toàn vẹn:} khoá ngoại, ràng buộc duy nhất và quy
        tắc xoá lan truyền (\code{CASCADE} / \code{SET NULL}) giữ dữ liệu
        nhất quán giữa từ vựng, bộ thẻ và tiến độ.
  \item \textbf{Chỉ mục:} đánh chỉ mục các cột truy vấn nóng (lemma, trạng
        thái ôn tập, thời điểm đến hạn) để bảo đảm hiệu năng khi dữ liệu lớn.
\end{itemize}

Bên cạnh PostgreSQL, hệ thống dùng \textbf{Redis} làm kho khoá--giá trị
trong bộ nhớ, đóng vai trò \textit{nơi lưu trữ hàng đợi} cho BullMQ
\cite{redis2024}. Việc truy cập cơ sở dữ liệu từ phía backend được thực hiện
qua \textbf{TypeORM} --- một thư viện ORM cho TypeScript/Node.js
\cite{typeorm2024} --- sẽ được trình bày ở mục~\ref{sec:cong-nghe-backend}.
Lược đồ cơ sở dữ liệu được quản lý bằng các \textit{migration} đặt trong thư
mục mã nguồn và chạy qua công cụ dòng lệnh của TypeORM, giúp tiến hoá lược
đồ một cách có kiểm soát và tái lập được giữa các môi trường.

% =====================================================================
\section{Công nghệ Frontend}
\label{sec:cong-nghe-frontend}

Ứng dụng khách (frontend) là một dự án độc lập, đóng vai trò bên tiêu thụ
API REST của backend. Frontend được xây dựng bằng \textbf{Next.js 16} theo
mô hình \textit{App Router} --- framework React hỗ trợ kết xuất phía máy chủ
(Server-Side Rendering --- SSR), \textit{thành phần phía máy chủ} (Server
Components) và \textit{hành động phía máy chủ} (Server Actions)
\cite{nextjs2024} --- trên nền thư viện giao diện \textbf{React~19}
\cite{react2024} và ngôn ngữ \textbf{TypeScript} (chế độ \textit{strict}).
Việc dùng chung TypeScript ở cả frontend và backend giúp thống nhất kiểu dữ
liệu của hợp đồng API giữa hai phía. Phần tạo kiểu giao diện dùng
\textbf{Tailwind CSS v4} --- framework CSS theo hướng tiện ích (utility-first),
cấu hình bằng CSS \cite{tailwind2024} --- kết hợp bộ thành phần
\textbf{shadcn/ui} dựng trên các primitive khả truy cập của \textbf{Base UI}.
Ràng buộc và kiểm tra dữ liệu biểu mẫu dùng \code{zod} \cite{zod2024}, được
dùng lại giữa kiểm tra phía máy khách và Server Action. Bảng~\ref{tab:frontend-libs}
liệt kê các thư viện frontend chính.

\begin{table}[h]
  \centering
  \small
  \renewcommand{\arraystretch}{1.3}
  \begin{tabular}{|p{4.2cm}|p{10.4cm}|}
    \hline
    \textbf{Công nghệ / thư viện} & \textbf{Vai trò} \\
    \hline
    Next.js 16 (App Router), React 19, TypeScript & Khung ứng dụng theo App Router; Server Components/Server Actions, định tuyến và kết xuất trang. \\
    \hline
    Tailwind CSS v4, shadcn/ui, Base UI & Hệ thống tạo kiểu tiện ích và bộ thành phần giao diện khả truy cập, tái sử dụng. \\
    \hline
    \code{zod} & Khai báo và kiểm tra ràng buộc dữ liệu (biểu mẫu, phản hồi API) phía máy khách. \\
    \hline
    \code{lucide-react}, \code{sonner} & Bộ biểu tượng và thông báo nổi (toast) cho phản hồi thao tác. \\
    \hline
    \code{next-themes}, \code{input-otp} & Hỗ trợ chế độ sáng/tối và ô nhập mã OTP (xác minh email). \\
    \hline
  \end{tabular}
  \caption{Các công nghệ và thư viện frontend chính}
  \label{tab:frontend-libs}
\end{table}

Mục này chỉ giới thiệu \textit{ngăn xếp công nghệ} frontend; quá trình hiện
thực hoá tầng giao diện (cấu trúc App Router, tầng truy cập API phía máy chủ,
quản lý phiên và động cơ phiên học phía máy khách) được trình bày ở
mục~\ref{sec:hien-thuc-frontend}, Chương~\ref{chuong:hien-thuc-hoa}. Trong mọi
trường hợp, frontend giao tiếp với backend hoàn toàn qua các API đã đặc tả,
không truy cập trực tiếp cơ sở dữ liệu.

% =====================================================================
\section{Công nghệ Backend}
\label{sec:cong-nghe-backend}

Backend được xây dựng bằng \textbf{NestJS}, một framework Node.js cho
TypeScript theo hướng mô-đun, dựa trên các nguyên lý \textit{tiêm phụ thuộc}
(dependency injection) và lập trình hướng khía cạnh \cite{nestjs2024}.
NestJS phù hợp với kiến trúc phân tầng đã nêu ở
mục~\ref{sec:kien-truc-tong-quan}: mỗi mô-đun đóng gói controller, service
và provider riêng; các \textit{guard}, \textit{pipe} và
\textit{interceptor} xử lý những mối quan tâm xuyên suốt (xác thực, kiểm tra
đầu vào, biến đổi phản hồi) một cách khai báo. Toàn bộ mã nguồn viết bằng
\textbf{TypeScript} ở chế độ \textit{strict}, giúp phát hiện lỗi kiểu ngay
khi biên dịch. Bảng~\ref{tab:backend-libs} liệt kê các thành phần backend
chính.

\begin{table}[h]
  \centering
  \small
  \renewcommand{\arraystretch}{1.3}
  \begin{tabular}{|p{4.6cm}|p{10.0cm}|}
    \hline
    \textbf{Công nghệ / thư viện} & \textbf{Vai trò} \\
    \hline
    NestJS 11 & Framework backend: tổ chức mô-đun, định tuyến, tiêm phụ thuộc, phiên bản hoá API theo URI. \\
    \hline
    TypeORM & Ánh xạ đối tượng--quan hệ tới PostgreSQL; quản lý thực thể và migration. \\
    \hline
    \code{class-validator}, \code{class-transformer} & Kiểm tra ràng buộc DTO và biến đổi dữ liệu vào/ra qua \code{ValidationPipe} toàn cục (loại bỏ trường lạ). \\
    \hline
    BullMQ + Redis & Hàng đợi công việc và worker xử lý nền (làm giàu từ vựng, sinh âm thanh, chấm điểm luyện tập). \\
    \hline
    \code{@nestjs/throttler} & Giới hạn tần suất (rate limiting) cho các điểm cuối nhạy cảm như đăng nhập. \\
    \hline
    \code{nodemailer} & Gửi email (mã xác minh) qua SMTP. \\
    \hline
    \code{pdf-parse}, \code{xlsx} & Trích xuất văn bản từ tệp PDF và bảng tính phục vụ nhập từ vựng hàng loạt. \\
    \hline
  \end{tabular}
  \caption{Các công nghệ và thư viện backend chính}
  \label{tab:backend-libs}
\end{table}

Phần xử lý bất đồng bộ được hiện thực bằng \textbf{BullMQ} --- thư viện hàng
đợi công việc chạy trên Redis \cite{bullmq2024}. Tiến trình API và tiến
trình worker chia sẻ chung cơ sở dữ liệu và hàng đợi nhưng được khởi chạy
độc lập (lệnh \code{start} và \code{start:worker}), phản ánh đúng mô hình
producer--consumer ở mục~\ref{sec:kien-truc-tong-quan}.

% =====================================================================
\section{Cơ chế xác thực và phân quyền}
\label{sec:xac-thuc-phan-quyen}

Phân hệ tài khoản áp dụng nhiều lớp kỹ thuật bảo mật, đáp ứng yêu cầu
NFR-01 trong Bảng~\ref{tab:nfr}.

\subsection{Băm mật khẩu}
\label{subsec:bcrypt}

Mật khẩu người dùng không bao giờ lưu ở dạng văn bản thuần mà được băm bằng
\textbf{bcrypt} --- hàm băm mật khẩu thích nghi, có \textit{muối} (salt) và
hệ số chi phí điều chỉnh được để chống tấn công vét cạn
\cite{provos1999bcrypt}. Khi đăng nhập, hệ thống so khớp mật khẩu nhập vào
với bản băm đã lưu thay vì giải mã.

\subsection{Xác thực bằng JSON Web Token}
\label{subsec:jwt}

Sau khi xác thực thành công, hệ thống cấp một cặp token theo chuẩn
\textbf{JSON Web Token (JWT)} \cite{rfc7519}: một \textit{access token} sống
ngắn để gọi API và một \textit{refresh token} sống dài để làm mới phiên.
JWT là chuỗi tự chứa, được ký số nên máy chủ có thể xác minh tính toàn vẹn
mà không cần tra cứu trạng thái phiên --- đúng tinh thần phi trạng thái của
REST (mục~\ref{subsec:rest-theory}). Thư viện \textbf{Passport} (qua
\code{passport-jwt}) đảm nhiệm việc trích và xác minh token trong mỗi yêu
cầu. Để có thể \textit{thu hồi}, refresh token được lưu dưới dạng băm trong
bảng \code{refresh_tokens} và bị vô hiệu khi người dùng đăng xuất.

\subsection{Phân quyền theo vai trò và quyền sở hữu}
\label{subsec:rbac}

Việc cấp quyền dựa trên hai cơ chế bổ trợ:

\begin{itemize}
  \item \textbf{Theo vai trò (role-based):} mỗi tài khoản mang vai trò
        \code{user} hoặc \code{admin}; một \textit{guard} kiểm tra vai trò
        để chặn các thao tác quản trị đối với người dùng thường.
  \item \textbf{Theo quyền sở hữu (ownership-based):} với tài nguyên cá
        nhân (từ vựng riêng, bộ thẻ, tiến độ), hệ thống kiểm tra người gọi
        có đúng là chủ sở hữu hay không trước khi cho phép truy cập.
\end{itemize}

\subsection{Đăng nhập qua nhà cung cấp bên thứ ba (OAuth/OIDC)}
\label{subsec:oauth}

Ngoài email/mật khẩu, hệ thống cho phép đăng nhập qua \textbf{Google},
\textbf{Apple} và \textbf{GitHub} bằng giao thức \textbf{OAuth~2.0} và lớp
nhận dạng \textbf{OpenID Connect} (OIDC) \cite{rfc6749}. Token danh tính do
nhà cung cấp cấp được hệ thống xác minh, sau đó liên kết với một bản ghi
trong bảng \code{user_identities}; nhờ vậy một người dùng có thể đăng nhập
bằng nhiều phương thức mà vẫn chung một tài khoản.

\subsection{Chấm điểm phiên học phi trạng thái bằng HMAC}
\label{subsec:hmac}

Một điểm đặc thù về bảo mật là cơ chế \textbf{chấm điểm phi trạng thái} cho
phiên học. Thay vì lưu trạng thái phiên trên máy chủ, mỗi câu hỏi khi sinh
ra được \textit{ký} bằng \textbf{HMAC} (Hash-based Message Authentication
Code) \cite{rfc2104} dựa trên một khoá bí mật của máy chủ. Khi người học nộp
đáp án, máy chủ xác minh lại chữ ký để bảo đảm câu hỏi không bị giả mạo hay
hết hạn, rồi mới suy lại đáp án đúng và chấm điểm. Cách làm này vừa giữ tính
phi trạng thái (thuận lợi cho mở rộng ngang) vừa chống gian lận.

% =====================================================================
\section{Công nghệ trí tuệ nhân tạo}
\label{sec:cong-nghe-ai}

Các chức năng thông minh của hệ thống dựa trên một số dịch vụ AI được tích
hợp qua tầng worker xử lý nền. Bảng~\ref{tab:ai-services} tổng hợp các dịch
vụ này.

\begin{table}[h]
  \centering
  \small
  \renewcommand{\arraystretch}{1.3}
  \begin{tabular}{|p{3.6cm}|p{4.0cm}|p{6.8cm}|}
    \hline
    \textbf{Dịch vụ} & \textbf{Loại} & \textbf{Vai trò trong hệ thống} \\
    \hline
    Gemma~3 (Google AI Studio) & Mô hình ngôn ngữ lớn (LLM) & Sinh nghĩa, ví dụ, bản dịch; chấm điểm câu luyện viết theo rubric kèm bậc CEFR. \\
    \hline
    Microsoft Edge TTS & Tổng hợp giọng nói (TTS) & Sinh âm thanh phát âm tự động cho từ vựng. \\
    \hline
    Vi dịch vụ chấm phát âm & Mô hình chấm điểm âm vị & Nhận bản ghi âm, trả điểm tổng và điểm từng phoneme. \\
    \hline
  \end{tabular}
  \caption{Các dịch vụ trí tuệ nhân tạo được tích hợp}
  \label{tab:ai-services}
\end{table}

\subsection{Sinh nội dung và chấm điểm bằng Gemma~3}
\label{subsec:gemma}

Hệ thống gọi mô hình \textbf{Gemma~3} qua nền tảng Google AI Studio
\cite{gemma2025technical}. Mô hình được sử dụng cho hai nhóm tác vụ:
(i)~\textit{làm giàu dữ liệu từ điển} --- từ một lemma, sinh các nghĩa, câu
ví dụ và bản dịch theo định dạng có cấu trúc; và (ii)~\textit{chấm điểm câu
luyện tập} --- đánh giá câu người học đặt theo thang điểm 0--100 kèm bậc
CEFR và nhận xét. Mọi lời gọi LLM đều là tác vụ nặng nên được đẩy sang
worker xử lý bất đồng bộ và bị giới hạn theo hạn mức ngày để bảo vệ khoá API
dùng chung.

\subsection{Tổng hợp giọng nói và chấm điểm phát âm}
\label{subsec:tts-pron}

Âm thanh phát âm của từ vựng được sinh tự động bằng \textbf{Microsoft Edge
TTS} (qua thư viện \code{msedge-tts}) rồi tải lên kho media. Chức năng luyện
phát âm được phục vụ bởi một \textbf{vi dịch vụ} riêng viết bằng Python
(FastAPI): backend đóng vai trò proxy mỏng, chuyển tiếp bản ghi âm của người
học tới vi dịch vụ để nhận điểm tổng và điểm từng âm vị, sau đó lưu lại lịch
sử luyện tập. Việc tách phần chấm phát âm thành dịch vụ độc lập giúp cô lập
các phụ thuộc nặng về xử lý tín hiệu/học máy khỏi tiến trình backend chính.

% =====================================================================
\section{Công cụ phát triển}
\label{sec:cong-cu-phat-trien}

Quá trình phát triển sử dụng bộ công cụ trong Bảng~\ref{tab:dev-tools} nhằm
bảo đảm chất lượng mã nguồn, khả năng tái lập môi trường và năng suất làm
việc.

\begin{table}[h]
  \centering
  \small
  \renewcommand{\arraystretch}{1.3}
  \begin{tabular}{|p{4.0cm}|p{10.6cm}|}
    \hline
    \textbf{Công cụ} & \textbf{Mục đích} \\
    \hline
    Git \& GitHub & Quản lý phiên bản mã nguồn và cộng tác qua pull request. \\
    \hline
    Docker \& Docker Compose & Đóng gói và khởi chạy nhất quán PostgreSQL, Redis, pgAdmin phục vụ phát triển và triển khai \cite{docker2024}. \\
    \hline
    NestJS CLI & Sinh khung mô-đun/controller/service và biên dịch, chạy ứng dụng. \\
    \hline
    TypeORM CLI & Sinh, chạy và hoàn tác migration của lược đồ cơ sở dữ liệu. \\
    \hline
    ESLint \& Prettier & Phân tích tĩnh và định dạng mã TypeScript theo quy ước thống nhất. \\
    \hline
    Jest \& Supertest & Kiểm thử đơn vị và kiểm thử tích hợp (end-to-end) cho API. \\
    \hline
    pgAdmin & Quản trị và truy vấn trực quan cơ sở dữ liệu PostgreSQL. \\
    \hline
    Overleaf (\LaTeX) & Biên soạn báo cáo đồ án theo mẫu của nhà trường. \\
    \hline
  \end{tabular}
  \caption{Các công cụ phát triển sử dụng trong đề tài}
  \label{tab:dev-tools}
\end{table}

Mã nguồn được quản lý bằng \textbf{Git} và lưu trên GitHub; mỗi tính năng
hoặc sửa lỗi được phát triển trên một nhánh riêng rồi hợp nhất qua pull
request. Môi trường phụ thuộc (PostgreSQL, Redis, pgAdmin) được mô tả bằng
\textbf{Docker Compose} \cite{docker2024}, giúp tái lập nhanh và đồng nhất
giữa các máy. Chất lượng mã được kiểm soát bằng kiểm tra kiểu của TypeScript,
\textbf{ESLint} và \textbf{Prettier}; phần kiểm thử dùng \textbf{Jest} cho
kiểm thử đơn vị và \textbf{Supertest} cho kiểm thử tích hợp các điểm cuối
API. Cuối cùng, báo cáo đồ án được biên soạn bằng \textbf{\LaTeX} trên
Overleaf theo đúng mẫu quy định.

% =====================================================================
\section{Kết chương}
\label{sec:ket-chuong-3}

Chương này đã trình bày các cơ sở lý thuyết --- ôn tập ngắt quãng và thuật
toán SM-2, khung trình độ CEFR, kiến trúc REST, ánh xạ đối tượng--quan hệ và
mô hình ngôn ngữ lớn --- cùng ngăn xếp công nghệ được lựa chọn cho từng tầng
của hệ thống, từ cơ sở dữ liệu, frontend, backend, cơ chế xác thực và phân
quyền, các dịch vụ trí tuệ nhân tạo đến bộ công cụ phát triển. Những lựa
chọn này trực tiếp phục vụ các quyết định thiết kế ở
Chương~\ref{chuong:phan-tich-thiet-ke} và là tiền đề cho quá trình hiện
thực hoá và kiểm thử được trình bày ở chương tiếp theo.

%  vào main.tex (sau Chương 2).
%
%  GHI CHÚ VỀ HÌNH VẼ: dùng macro \figph{<tên-file>}{<tỉ-lệ-rộng>}{<chú-thích>}{<nhãn>}
%  (đã khai báo ở Chương 2, khai báo lại bằng \providecommand để an toàn nếu
%  biên dịch độc lập). Macro tự chèn ảnh thật nếu tệp tồn tại trong images/,
%  ngược lại hiển thị khung giữ chỗ để báo cáo vẫn biên dịch được.
%
%  KHÔNG hardcode số chương/mục/bảng/hình/công thức — luôn dùng \label + \ref / \cite.
% =====================================================================

% --- Macro tiện ích (idempotent — không định nghĩa lại nếu đã có) ------
\providecommand{\code}[1]{\texttt{\detokenize{#1}}} % in định danh snake_case an toàn
\providecommand{\figph}[4]{%
  \begin{figure}[htbp]
    \centering
    \IfFileExists{images/#1}%
      {\includegraphics[width=#2\textwidth,height=0.82\textheight,keepaspectratio]{#1}}%
      {\setlength{\fboxsep}{10pt}\fbox{\begin{minipage}[c][3.6cm][c]{#2\textwidth}\centering\itshape Sơ đồ minh hoạ --- chèn tệp ảnh \code{images/#1} vào thư mục \code{images/}.\end{minipage}}}%
    \caption{#3}%
    \label{#4}%
  \end{figure}}

\chapter{CƠ SỞ LÝ THUYẾT VÀ CÔNG NGHỆ SỬ DỤNG}
\label{chuong:co-so-ly-thuyet}

Chương~\ref{chuong:phan-tich-thiet-ke} đã trình bày kết quả phân tích và
thiết kế hệ thống. Chương này hệ thống hoá những \textit{cơ sở lý thuyết}
và \textit{công nghệ} làm nền tảng cho các quyết định thiết kế và hiện thực
hoá đó, qua đó lý giải vì sao mỗi công nghệ được lựa chọn. Trước hết,
mục~\ref{sec:co-so-ly-thuyet} trình bày các cơ sở lý thuyết cốt lõi: ôn tập
ngắt quãng và thuật toán SM-2, khung trình độ CEFR, kiến trúc REST, ánh xạ
đối tượng--quan hệ và mô hình ngôn ngữ lớn. Tiếp theo,
mục~\ref{sec:tong-quan-cong-nghe} tổng quan ngăn xếp công nghệ và
mục~\ref{sec:kien-truc-tong-quan} phân tích các kiểu kiến trúc được áp dụng.
Các mục~\ref{sec:cong-nghe-csdl}--\ref{sec:cong-cu-phat-trien} đi sâu lần
lượt vào công nghệ cơ sở dữ liệu, frontend, backend, cơ chế xác thực và
phân quyền, công nghệ trí tuệ nhân tạo và bộ công cụ phát triển.

% =====================================================================
\section{Cơ sở lý thuyết}
\label{sec:co-so-ly-thuyet}

\subsection{Ôn tập ngắt quãng và đường cong lãng quên}
\label{subsec:srs-theory}

Cơ sở khoa học của động cơ học trong hệ thống là hiện tượng
\textit{đường cong lãng quên} (forgetting curve) do Ebbinghaus phát hiện:
khả năng ghi nhớ một thông tin mới suy giảm theo hàm mũ theo thời gian nếu
không được củng cố \cite{ebbinghaus1913memory}. Một cách mô hình hoá đơn
giản, mức ghi nhớ \( R \) còn lại sau khoảng thời gian \( t \) kể từ lần
học gần nhất được biểu diễn qua công thức~\eqref{eq:forgetting}:

\begin{equation}
  \label{eq:forgetting}
  R(t) = e^{-\,t/S},
\end{equation}

trong đó \( S \) là \textit{độ bền trí nhớ} (memory stability) --- càng ôn
tập đúng lúc, \( S \) càng lớn và thông tin được giữ lâu hơn. Hệ quả thực
tiễn là: thay vì học nhồi nhét, người học nên ôn lại mỗi từ tại những
\textit{thời điểm đến hạn} được giãn dần. Các nghiên cứu về hiệu ứng giãn
cách (spacing effect) khẳng định ôn tập ngắt quãng cải thiện rõ rệt khả
năng ghi nhớ dài hạn so với học dồn \cite{cepeda2006distributed}. Đây là
nền tảng lý thuyết cho \textit{hệ thống ôn tập ngắt quãng} (Spaced
Repetition System --- SRS) mà hệ thống hiện thực hoá.

\subsection{Thuật toán lập lịch ôn tập SM-2}
\label{subsec:sm2-theory}

Để quyết định \textit{khi nào} ôn lại mỗi từ, hệ thống áp dụng thuật toán
\textbf{SM-2} có nguồn gốc từ dự án SuperMemo \cite{wozniak1990optimization}.
Mỗi cặp (người dùng, từ) được mô tả bằng ba đại lượng: hệ số dễ nhớ
\textit{easiness factor} \( EF \) (khởi tạo \( 2{,}5 \)), số lần ôn đúng
liên tiếp \( n \) và khoảng cách ôn tập \( I \) (tính theo ngày). Sau mỗi
lượt ôn, người học (ở đây là kết quả trả lời trong phiên) được quy về một
\textit{điểm chất lượng} \( q \in \{0,1,\dots,5\} \).

Hệ số \( EF \) được cập nhật theo công thức~\eqref{eq:sm2-ef}, có chặn dưới
\( 1{,}3 \) để khoảng ôn tập không co lại quá nhanh:

\begin{equation}
  \label{eq:sm2-ef}
  EF' = EF + \big(0{,}1 - (5-q)\,(0{,}08 + (5-q)\cdot 0{,}02)\big),
  \qquad EF' \geq 1{,}3 .
\end{equation}

Khoảng cách ôn tập kế tiếp được tính theo công thức~\eqref{eq:sm2-interval}:

\begin{equation}
  \label{eq:sm2-interval}
  I(n) =
  \begin{cases}
    1, & n = 1,\\[2pt]
    6, & n = 2,\\[2pt]
    \big\lceil I(n-1)\times EF' \big\rceil, & n > 2 .
  \end{cases}
\end{equation}

Nếu người học trả lời sai (\( q < 3 \)), bộ đếm \( n \) được đặt lại về 0 và
từ phải học lại từ đầu; ngược lại \( n \) tăng và khoảng ôn tập giãn dần
theo \( EF' \). Trong hệ thống, SM-2 được \textbf{mở rộng bằng các bước học}
(learning steps): một từ mới phải vượt qua một chuỗi câu hỏi từ dễ đến khó
trong phiên trước khi được giao cho SM-2 lập lịch dài hạn; điểm chất lượng
\( q \) được suy ra từ độ chính xác ở bước cuối của mỗi từ. Trạng thái thẻ
nhờ đó chuyển dịch \code{new} \( \to \) \code{learning} \( \to \)
\code{review} \( \to \) \code{mastered}, phù hợp với biểu đồ hoạt động đã
trình bày ở Chương~\ref{chuong:phan-tich-thiet-ke}.

\subsection{Khung tham chiếu trình độ ngôn ngữ CEFR}
\label{subsec:cefr-theory}

Hệ thống chuẩn hoá độ khó của từ vựng và trình độ người học theo
\textbf{Khung tham chiếu trình độ ngôn ngữ chung châu Âu} (Common European
Framework of Reference for Languages --- CEFR) do Hội đồng châu Âu ban hành
\cite{coe2001cefr}. CEFR chia năng lực ngôn ngữ thành sáu bậc từ
\( A1, A2 \) (sơ cấp), \( B1, B2 \) (trung cấp) đến \( C1, C2 \) (cao cấp).
Mỗi từ trong kho được gắn một bậc CEFR và mỗi người học khai báo trình độ
của mình; nhờ đó hệ thống lựa chọn từ và gợi ý bộ thẻ phù hợp với năng lực
hiện tại, đồng thời mô hình ngôn ngữ chấm điểm bậc CEFR mà một câu luyện
tập thể hiện.

\subsection{Kiến trúc REST và dịch vụ web}
\label{subsec:rest-theory}

Giao diện lập trình ứng dụng (API) của hệ thống tuân theo phong cách kiến
trúc \textbf{REST} (Representational State Transfer) do Fielding đề xuất
\cite{fielding2000rest}. Các ràng buộc cốt lõi của REST mà hệ thống áp dụng
gồm: mô hình \textit{client--server} tách bạch máy khách và máy chủ; tính
\textit{phi trạng thái} (stateless) --- mỗi yêu cầu tự mang đủ thông tin xác
thực (token), máy chủ không lưu trạng thái phiên; định danh tài nguyên bằng
\textit{URI} và thao tác qua các phương thức HTTP chuẩn (GET, POST, PATCH,
DELETE) với ngữ nghĩa thống nhất; và trao đổi dữ liệu dưới dạng
\textbf{JSON}. Tính phi trạng thái là điều kiện để mở rộng tầng API theo
chiều ngang (đáp ứng yêu cầu NFR-03 trong
Bảng~\ref{tab:nfr}). API còn được \textit{phiên bản hoá theo URI}
(tiền tố \code{/v1}) nhằm tiến hoá hợp đồng mà không phá vỡ máy khách cũ.

\subsection{Ánh xạ đối tượng--quan hệ (ORM)}
\label{subsec:orm-theory}

Để thu hẹp khoảng cách giữa mô hình hướng đối tượng trong mã nguồn và mô
hình quan hệ trong cơ sở dữ liệu, hệ thống dùng kỹ thuật \textbf{ánh xạ đối
tượng--quan hệ} (Object--Relational Mapping --- ORM). Mỗi bảng được khai
báo dưới dạng một lớp \textit{thực thể} (entity), mỗi dòng tương ứng một đối
tượng và các quan hệ khoá ngoại được biểu diễn bằng liên kết giữa các thực
thể. ORM giúp lập trình viên thao tác dữ liệu bằng mã có kiểu (typed) thay
vì viết câu lệnh SQL thủ công, đồng thời quản lý tiến hoá lược đồ bằng cơ
chế \textit{migration} có thể phiên bản hoá. Đánh đổi của ORM là một lớp
trừu tượng phụ và nguy cơ truy vấn kém hiệu quả nếu lạm dụng; hệ thống khắc
phục bằng cách đánh chỉ mục các truy vấn nóng và truy vấn tường minh khi cần.

\subsection{Mô hình ngôn ngữ lớn và sinh nội dung}
\label{subsec:llm-theory}

Các chức năng hỗ trợ bằng AI (làm giàu dữ liệu từ điển, chấm điểm câu luyện
tập) dựa trên \textbf{mô hình ngôn ngữ lớn} (Large Language Model --- LLM).
LLM là mạng nơ-ron kiến trúc \textit{Transformer} được huấn luyện trên khối
dữ liệu văn bản khổng lồ để dự đoán token kế tiếp, nhờ đó có khả năng sinh
văn bản mạch lạc, dịch và đánh giá ngôn ngữ theo chỉ dẫn (prompt). Hệ thống
sử dụng họ mô hình mở \textbf{Gemma} của Google DeepMind
\cite{gemma2025technical}. Trong nghiệp vụ, LLM được giao vai trò có kiểm
soát: sinh nghĩa, ví dụ, bản dịch theo định dạng yêu cầu và chấm điểm câu
theo \textit{rubric}; kết quả do AI sinh ra đối với kho hệ thống luôn phải
qua bước \textit{duyệt} của quản trị viên trước khi xuất bản, nhằm kiểm soát
chất lượng và hạn chế lỗi "ảo giác" (hallucination) của mô hình.

% =====================================================================
\section{Tổng quan công nghệ sử dụng}
\label{sec:tong-quan-cong-nghe}

Trên cơ sở lý thuyết ở mục~\ref{sec:co-so-ly-thuyet}, hệ thống lựa chọn ngăn
xếp công nghệ được tổ chức theo nhiều tầng. Bảng~\ref{tab:tech-overview}
tóm tắt công nghệ chủ chốt của từng tầng cùng vai trò của nó;
Hình~\ref{fig:tech-stack} minh hoạ ngăn xếp này một cách trực quan. Các mục
tiếp theo phân tích chi tiết từng nhóm công nghệ.

\begin{table}[h]
  \centering
  \small
  \renewcommand{\arraystretch}{1.3}
  \begin{tabular}{|p{3.0cm}|p{5.0cm}|p{6.4cm}|}
    \hline
    \textbf{Tầng} & \textbf{Công nghệ chính} & \textbf{Vai trò} \\
    \hline
    Frontend & Next.js, React, TypeScript, Tailwind CSS, shadcn/ui & Ứng dụng khách tiêu thụ API, dựng giao diện người dùng. \\
    \hline
    Backend & NestJS 11, TypeScript, TypeORM & Cung cấp API REST, chứa nghiệp vụ và truy cập dữ liệu. \\
    \hline
    Cơ sở dữ liệu & PostgreSQL & Lưu trữ dữ liệu quan hệ có ràng buộc toàn vẹn. \\
    \hline
    Hàng đợi \& nền & Redis, BullMQ & Hàng đợi công việc và worker xử lý bất đồng bộ. \\
    \hline
    Xác thực & JWT, Passport, bcrypt, OAuth/OIDC & Xác thực, phân quyền và bảo vệ tài khoản. \\
    \hline
    Trí tuệ nhân tạo & Gemma~3 (Google AI Studio), Edge TTS, vi dịch vụ chấm phát âm & Sinh nội dung, sinh âm thanh, chấm điểm luyện tập. \\
    \hline
    Hạ tầng phụ trợ & Cloudinary, SMTP, Docker & Lưu trữ media, gửi email, đóng gói triển khai. \\
    \hline
  \end{tabular}
  \caption{Tổng quan ngăn xếp công nghệ của hệ thống}
  \label{tab:tech-overview}
\end{table}

% --- PlantUML gợi ý cho Hình ngăn xếp công nghệ ----------------------
% @startuml
% rectangle "FRONTEND\nNext.js + React + TypeScript\nTailwind CSS + shadcn/ui" as FE
% rectangle "BACKEND (API)\nNestJS 11 + TypeScript\nController - Service - Repository (TypeORM)" as BE
% rectangle "XỬ LÝ NỀN\nWorker (BullMQ) + Redis" as WK
% database "PostgreSQL" as DB
% cloud "DỊCH VỤ AI / NGOÀI\nGemma 3, Edge TTS, Cloudinary,\nDV chấm phát âm, OAuth, SMTP" as EXT
% FE --> BE : HTTPS / JSON (REST)
% BE --> DB
% BE --> WK : enqueue
% WK --> DB
% WK --> EXT
% BE --> EXT
% @enduml
\figph{tech-stack.png}{0.8}{Ngăn xếp công nghệ của hệ thống theo tầng}{fig:tech-stack}

% =====================================================================
\section{Tổng quan kiến trúc hệ thống}
\label{sec:kien-truc-tong-quan}

Mục này điểm lại các \textit{kiểu kiến trúc} (architectural styles) làm nền
cho thiết kế đã trình bày chi tiết ở mục kiến trúc của
Chương~\ref{chuong:phan-tich-thiet-ke} (xem Hình~\ref{fig:kien-truc}).

\begin{itemize}
  \item \textbf{Monolith mô-đun hoá (modular monolith):} toàn bộ backend là
        một tiến trình triển khai duy nhất nhưng được chia thành các mô-đun
        nghiệp vụ tách bạch (xem Bảng~\ref{tab:modules}). So với kiến trúc
        vi dịch vụ (microservices), cách này đơn giản hơn khi triển khai và
        vận hành ở quy mô đồ án, đồng thời vẫn giữ ranh giới mô-đun rõ ràng
        để có thể tách dịch vụ về sau.
  \item \textbf{Kiến trúc phân tầng (layered):} trong mỗi mô-đun, mã nguồn
        chia theo các tầng \textit{Controller} (tiếp nhận HTTP, kiểm tra
        DTO) \( \to \) \textit{Service} (nghiệp vụ) \( \to \)
        \textit{Repository} (truy cập dữ liệu qua ORM). Việc tách tầng giúp
        kiểm thử và bảo trì độc lập từng mối quan tâm.
  \item \textbf{Client--server phi trạng thái (REST):} máy khách
        (Next.js) và máy chủ giao tiếp qua HTTP/JSON theo REST; máy chủ phi
        trạng thái nên có thể nhân bản để mở rộng theo chiều ngang.
  \item \textbf{Xử lý bất đồng bộ kiểu nhà sản xuất--người tiêu thụ
        (producer--consumer):} tiến trình API đóng vai trò nhà sản xuất, đẩy
        công việc nặng vào hàng đợi; tiến trình worker là người tiêu thụ,
        xử lý nền. Mẫu này giữ cho luồng yêu cầu chính luôn nhanh và ổn định.
\end{itemize}

Sự kết hợp bốn kiểu kiến trúc trên cho phép hệ thống vừa đơn giản khi triển
khai, vừa đáp ứng các yêu cầu phi chức năng về hiệu năng, khả năng mở rộng
và khả năng bảo trì trong Bảng~\ref{tab:nfr}.

% =====================================================================
\section{Công nghệ cơ sở dữ liệu}
\label{sec:cong-nghe-csdl}

Hệ thống dùng hệ quản trị cơ sở dữ liệu quan hệ \textbf{PostgreSQL}
\cite{postgresql2024}. PostgreSQL được chọn vì là mã nguồn mở, trưởng thành,
tuân thủ chuẩn SQL và hỗ trợ tốt các đặc tính mà thiết kế dữ liệu
(mục thiết kế CSDL của Chương~\ref{chuong:phan-tich-thiet-ke}) yêu cầu:

\begin{itemize}
  \item \textbf{Giao dịch ACID:} bảo đảm tính nguyên tử cho các thao tác
        ghi nhiều bảng (ví dụ tạo phiên học, cập nhật tiến độ và ghi nhật
        ký), đáp ứng yêu cầu độ tin cậy.
  \item \textbf{Khoá chính \code{uuid}:} dùng kiểu định danh duy nhất toàn
        cục cho mọi bảng, thuận lợi cho phân tán và tránh lộ số lượng bản ghi.
  \item \textbf{Kiểu \code{timestamptz}:} lưu mốc thời gian kèm múi giờ,
        bảo đảm nhất quán cho thống kê chuỗi ngày học và lịch ôn tập.
  \item \textbf{Ràng buộc toàn vẹn:} khoá ngoại, ràng buộc duy nhất và quy
        tắc xoá lan truyền (\code{CASCADE} / \code{SET NULL}) giữ dữ liệu
        nhất quán giữa từ vựng, bộ thẻ và tiến độ.
  \item \textbf{Chỉ mục:} đánh chỉ mục các cột truy vấn nóng (lemma, trạng
        thái ôn tập, thời điểm đến hạn) để bảo đảm hiệu năng khi dữ liệu lớn.
\end{itemize}

Bên cạnh PostgreSQL, hệ thống dùng \textbf{Redis} làm kho khoá--giá trị
trong bộ nhớ, đóng vai trò \textit{nơi lưu trữ hàng đợi} cho BullMQ
\cite{redis2024}. Việc truy cập cơ sở dữ liệu từ phía backend được thực hiện
qua \textbf{TypeORM} --- một thư viện ORM cho TypeScript/Node.js
\cite{typeorm2024} --- sẽ được trình bày ở mục~\ref{sec:cong-nghe-backend}.
Lược đồ cơ sở dữ liệu được quản lý bằng các \textit{migration} đặt trong thư
mục mã nguồn và chạy qua công cụ dòng lệnh của TypeORM, giúp tiến hoá lược
đồ một cách có kiểm soát và tái lập được giữa các môi trường.

% =====================================================================
\section{Công nghệ Frontend}
\label{sec:cong-nghe-frontend}

Ứng dụng khách (frontend) là một dự án độc lập, đóng vai trò bên tiêu thụ
API REST của backend. Frontend được xây dựng bằng \textbf{Next.js 16} theo
mô hình \textit{App Router} --- framework React hỗ trợ kết xuất phía máy chủ
(Server-Side Rendering --- SSR), \textit{thành phần phía máy chủ} (Server
Components) và \textit{hành động phía máy chủ} (Server Actions)
\cite{nextjs2024} --- trên nền thư viện giao diện \textbf{React~19}
\cite{react2024} và ngôn ngữ \textbf{TypeScript} (chế độ \textit{strict}).
Việc dùng chung TypeScript ở cả frontend và backend giúp thống nhất kiểu dữ
liệu của hợp đồng API giữa hai phía. Phần tạo kiểu giao diện dùng
\textbf{Tailwind CSS v4} --- framework CSS theo hướng tiện ích (utility-first),
cấu hình bằng CSS \cite{tailwind2024} --- kết hợp bộ thành phần
\textbf{shadcn/ui} dựng trên các primitive khả truy cập của \textbf{Base UI}.
Ràng buộc và kiểm tra dữ liệu biểu mẫu dùng \code{zod} \cite{zod2024}, được
dùng lại giữa kiểm tra phía máy khách và Server Action. Bảng~\ref{tab:frontend-libs}
liệt kê các thư viện frontend chính.

\begin{table}[h]
  \centering
  \small
  \renewcommand{\arraystretch}{1.3}
  \begin{tabular}{|p{4.2cm}|p{10.4cm}|}
    \hline
    \textbf{Công nghệ / thư viện} & \textbf{Vai trò} \\
    \hline
    Next.js 16 (App Router), React 19, TypeScript & Khung ứng dụng theo App Router; Server Components/Server Actions, định tuyến và kết xuất trang. \\
    \hline
    Tailwind CSS v4, shadcn/ui, Base UI & Hệ thống tạo kiểu tiện ích và bộ thành phần giao diện khả truy cập, tái sử dụng. \\
    \hline
    \code{zod} & Khai báo và kiểm tra ràng buộc dữ liệu (biểu mẫu, phản hồi API) phía máy khách. \\
    \hline
    \code{lucide-react}, \code{sonner} & Bộ biểu tượng và thông báo nổi (toast) cho phản hồi thao tác. \\
    \hline
    \code{next-themes}, \code{input-otp} & Hỗ trợ chế độ sáng/tối và ô nhập mã OTP (xác minh email). \\
    \hline
  \end{tabular}
  \caption{Các công nghệ và thư viện frontend chính}
  \label{tab:frontend-libs}
\end{table}

Mục này chỉ giới thiệu \textit{ngăn xếp công nghệ} frontend; quá trình hiện
thực hoá tầng giao diện (cấu trúc App Router, tầng truy cập API phía máy chủ,
quản lý phiên và động cơ phiên học phía máy khách) được trình bày ở
mục~\ref{sec:hien-thuc-frontend}, Chương~\ref{chuong:hien-thuc-hoa}. Trong mọi
trường hợp, frontend giao tiếp với backend hoàn toàn qua các API đã đặc tả,
không truy cập trực tiếp cơ sở dữ liệu.

% =====================================================================
\section{Công nghệ Backend}
\label{sec:cong-nghe-backend}

Backend được xây dựng bằng \textbf{NestJS}, một framework Node.js cho
TypeScript theo hướng mô-đun, dựa trên các nguyên lý \textit{tiêm phụ thuộc}
(dependency injection) và lập trình hướng khía cạnh \cite{nestjs2024}.
NestJS phù hợp với kiến trúc phân tầng đã nêu ở
mục~\ref{sec:kien-truc-tong-quan}: mỗi mô-đun đóng gói controller, service
và provider riêng; các \textit{guard}, \textit{pipe} và
\textit{interceptor} xử lý những mối quan tâm xuyên suốt (xác thực, kiểm tra
đầu vào, biến đổi phản hồi) một cách khai báo. Toàn bộ mã nguồn viết bằng
\textbf{TypeScript} ở chế độ \textit{strict}, giúp phát hiện lỗi kiểu ngay
khi biên dịch. Bảng~\ref{tab:backend-libs} liệt kê các thành phần backend
chính.

\begin{table}[h]
  \centering
  \small
  \renewcommand{\arraystretch}{1.3}
  \begin{tabular}{|p{4.6cm}|p{10.0cm}|}
    \hline
    \textbf{Công nghệ / thư viện} & \textbf{Vai trò} \\
    \hline
    NestJS 11 & Framework backend: tổ chức mô-đun, định tuyến, tiêm phụ thuộc, phiên bản hoá API theo URI. \\
    \hline
    TypeORM & Ánh xạ đối tượng--quan hệ tới PostgreSQL; quản lý thực thể và migration. \\
    \hline
    \code{class-validator}, \code{class-transformer} & Kiểm tra ràng buộc DTO và biến đổi dữ liệu vào/ra qua \code{ValidationPipe} toàn cục (loại bỏ trường lạ). \\
    \hline
    BullMQ + Redis & Hàng đợi công việc và worker xử lý nền (làm giàu từ vựng, sinh âm thanh, chấm điểm luyện tập). \\
    \hline
    \code{@nestjs/throttler} & Giới hạn tần suất (rate limiting) cho các điểm cuối nhạy cảm như đăng nhập. \\
    \hline
    \code{nodemailer} & Gửi email (mã xác minh) qua SMTP. \\
    \hline
    \code{pdf-parse}, \code{xlsx} & Trích xuất văn bản từ tệp PDF và bảng tính phục vụ nhập từ vựng hàng loạt. \\
    \hline
  \end{tabular}
  \caption{Các công nghệ và thư viện backend chính}
  \label{tab:backend-libs}
\end{table}

Phần xử lý bất đồng bộ được hiện thực bằng \textbf{BullMQ} --- thư viện hàng
đợi công việc chạy trên Redis \cite{bullmq2024}. Tiến trình API và tiến
trình worker chia sẻ chung cơ sở dữ liệu và hàng đợi nhưng được khởi chạy
độc lập (lệnh \code{start} và \code{start:worker}), phản ánh đúng mô hình
producer--consumer ở mục~\ref{sec:kien-truc-tong-quan}.

% =====================================================================
\section{Cơ chế xác thực và phân quyền}
\label{sec:xac-thuc-phan-quyen}

Phân hệ tài khoản áp dụng nhiều lớp kỹ thuật bảo mật, đáp ứng yêu cầu
NFR-01 trong Bảng~\ref{tab:nfr}.

\subsection{Băm mật khẩu}
\label{subsec:bcrypt}

Mật khẩu người dùng không bao giờ lưu ở dạng văn bản thuần mà được băm bằng
\textbf{bcrypt} --- hàm băm mật khẩu thích nghi, có \textit{muối} (salt) và
hệ số chi phí điều chỉnh được để chống tấn công vét cạn
\cite{provos1999bcrypt}. Khi đăng nhập, hệ thống so khớp mật khẩu nhập vào
với bản băm đã lưu thay vì giải mã.

\subsection{Xác thực bằng JSON Web Token}
\label{subsec:jwt}

Sau khi xác thực thành công, hệ thống cấp một cặp token theo chuẩn
\textbf{JSON Web Token (JWT)} \cite{rfc7519}: một \textit{access token} sống
ngắn để gọi API và một \textit{refresh token} sống dài để làm mới phiên.
JWT là chuỗi tự chứa, được ký số nên máy chủ có thể xác minh tính toàn vẹn
mà không cần tra cứu trạng thái phiên --- đúng tinh thần phi trạng thái của
REST (mục~\ref{subsec:rest-theory}). Thư viện \textbf{Passport} (qua
\code{passport-jwt}) đảm nhiệm việc trích và xác minh token trong mỗi yêu
cầu. Để có thể \textit{thu hồi}, refresh token được lưu dưới dạng băm trong
bảng \code{refresh_tokens} và bị vô hiệu khi người dùng đăng xuất.

\subsection{Phân quyền theo vai trò và quyền sở hữu}
\label{subsec:rbac}

Việc cấp quyền dựa trên hai cơ chế bổ trợ:

\begin{itemize}
  \item \textbf{Theo vai trò (role-based):} mỗi tài khoản mang vai trò
        \code{user} hoặc \code{admin}; một \textit{guard} kiểm tra vai trò
        để chặn các thao tác quản trị đối với người dùng thường.
  \item \textbf{Theo quyền sở hữu (ownership-based):} với tài nguyên cá
        nhân (từ vựng riêng, bộ thẻ, tiến độ), hệ thống kiểm tra người gọi
        có đúng là chủ sở hữu hay không trước khi cho phép truy cập.
\end{itemize}

\subsection{Đăng nhập qua nhà cung cấp bên thứ ba (OAuth/OIDC)}
\label{subsec:oauth}

Ngoài email/mật khẩu, hệ thống cho phép đăng nhập qua \textbf{Google},
\textbf{Apple} và \textbf{GitHub} bằng giao thức \textbf{OAuth~2.0} và lớp
nhận dạng \textbf{OpenID Connect} (OIDC) \cite{rfc6749}. Token danh tính do
nhà cung cấp cấp được hệ thống xác minh, sau đó liên kết với một bản ghi
trong bảng \code{user_identities}; nhờ vậy một người dùng có thể đăng nhập
bằng nhiều phương thức mà vẫn chung một tài khoản.

\subsection{Chấm điểm phiên học phi trạng thái bằng HMAC}
\label{subsec:hmac}

Một điểm đặc thù về bảo mật là cơ chế \textbf{chấm điểm phi trạng thái} cho
phiên học. Thay vì lưu trạng thái phiên trên máy chủ, mỗi câu hỏi khi sinh
ra được \textit{ký} bằng \textbf{HMAC} (Hash-based Message Authentication
Code) \cite{rfc2104} dựa trên một khoá bí mật của máy chủ. Khi người học nộp
đáp án, máy chủ xác minh lại chữ ký để bảo đảm câu hỏi không bị giả mạo hay
hết hạn, rồi mới suy lại đáp án đúng và chấm điểm. Cách làm này vừa giữ tính
phi trạng thái (thuận lợi cho mở rộng ngang) vừa chống gian lận.

% =====================================================================
\section{Công nghệ trí tuệ nhân tạo}
\label{sec:cong-nghe-ai}

Các chức năng thông minh của hệ thống dựa trên một số dịch vụ AI được tích
hợp qua tầng worker xử lý nền. Bảng~\ref{tab:ai-services} tổng hợp các dịch
vụ này.

\begin{table}[h]
  \centering
  \small
  \renewcommand{\arraystretch}{1.3}
  \begin{tabular}{|p{3.6cm}|p{4.0cm}|p{6.8cm}|}
    \hline
    \textbf{Dịch vụ} & \textbf{Loại} & \textbf{Vai trò trong hệ thống} \\
    \hline
    Gemma~3 (Google AI Studio) & Mô hình ngôn ngữ lớn (LLM) & Sinh nghĩa, ví dụ, bản dịch; chấm điểm câu luyện viết theo rubric kèm bậc CEFR. \\
    \hline
    Microsoft Edge TTS & Tổng hợp giọng nói (TTS) & Sinh âm thanh phát âm tự động cho từ vựng. \\
    \hline
    Vi dịch vụ chấm phát âm & Mô hình chấm điểm âm vị & Nhận bản ghi âm, trả điểm tổng và điểm từng phoneme. \\
    \hline
  \end{tabular}
  \caption{Các dịch vụ trí tuệ nhân tạo được tích hợp}
  \label{tab:ai-services}
\end{table}

\subsection{Sinh nội dung và chấm điểm bằng Gemma~3}
\label{subsec:gemma}

Hệ thống gọi mô hình \textbf{Gemma~3} qua nền tảng Google AI Studio
\cite{gemma2025technical}. Mô hình được sử dụng cho hai nhóm tác vụ:
(i)~\textit{làm giàu dữ liệu từ điển} --- từ một lemma, sinh các nghĩa, câu
ví dụ và bản dịch theo định dạng có cấu trúc; và (ii)~\textit{chấm điểm câu
luyện tập} --- đánh giá câu người học đặt theo thang điểm 0--100 kèm bậc
CEFR và nhận xét. Mọi lời gọi LLM đều là tác vụ nặng nên được đẩy sang
worker xử lý bất đồng bộ và bị giới hạn theo hạn mức ngày để bảo vệ khoá API
dùng chung.

\subsection{Tổng hợp giọng nói và chấm điểm phát âm}
\label{subsec:tts-pron}

Âm thanh phát âm của từ vựng được sinh tự động bằng \textbf{Microsoft Edge
TTS} (qua thư viện \code{msedge-tts}) rồi tải lên kho media. Chức năng luyện
phát âm được phục vụ bởi một \textbf{vi dịch vụ} riêng viết bằng Python
(FastAPI): backend đóng vai trò proxy mỏng, chuyển tiếp bản ghi âm của người
học tới vi dịch vụ để nhận điểm tổng và điểm từng âm vị, sau đó lưu lại lịch
sử luyện tập. Việc tách phần chấm phát âm thành dịch vụ độc lập giúp cô lập
các phụ thuộc nặng về xử lý tín hiệu/học máy khỏi tiến trình backend chính.

% =====================================================================
\section{Công cụ phát triển}
\label{sec:cong-cu-phat-trien}

Quá trình phát triển sử dụng bộ công cụ trong Bảng~\ref{tab:dev-tools} nhằm
bảo đảm chất lượng mã nguồn, khả năng tái lập môi trường và năng suất làm
việc.

\begin{table}[h]
  \centering
  \small
  \renewcommand{\arraystretch}{1.3}
  \begin{tabular}{|p{4.0cm}|p{10.6cm}|}
    \hline
    \textbf{Công cụ} & \textbf{Mục đích} \\
    \hline
    Git \& GitHub & Quản lý phiên bản mã nguồn và cộng tác qua pull request. \\
    \hline
    Docker \& Docker Compose & Đóng gói và khởi chạy nhất quán PostgreSQL, Redis, pgAdmin phục vụ phát triển và triển khai \cite{docker2024}. \\
    \hline
    NestJS CLI & Sinh khung mô-đun/controller/service và biên dịch, chạy ứng dụng. \\
    \hline
    TypeORM CLI & Sinh, chạy và hoàn tác migration của lược đồ cơ sở dữ liệu. \\
    \hline
    ESLint \& Prettier & Phân tích tĩnh và định dạng mã TypeScript theo quy ước thống nhất. \\
    \hline
    Jest \& Supertest & Kiểm thử đơn vị và kiểm thử tích hợp (end-to-end) cho API. \\
    \hline
    pgAdmin & Quản trị và truy vấn trực quan cơ sở dữ liệu PostgreSQL. \\
    \hline
    Overleaf (\LaTeX) & Biên soạn báo cáo đồ án theo mẫu của nhà trường. \\
    \hline
  \end{tabular}
  \caption{Các công cụ phát triển sử dụng trong đề tài}
  \label{tab:dev-tools}
\end{table}

Mã nguồn được quản lý bằng \textbf{Git} và lưu trên GitHub; mỗi tính năng
hoặc sửa lỗi được phát triển trên một nhánh riêng rồi hợp nhất qua pull
request. Môi trường phụ thuộc (PostgreSQL, Redis, pgAdmin) được mô tả bằng
\textbf{Docker Compose} \cite{docker2024}, giúp tái lập nhanh và đồng nhất
giữa các máy. Chất lượng mã được kiểm soát bằng kiểm tra kiểu của TypeScript,
\textbf{ESLint} và \textbf{Prettier}; phần kiểm thử dùng \textbf{Jest} cho
kiểm thử đơn vị và \textbf{Supertest} cho kiểm thử tích hợp các điểm cuối
API. Cuối cùng, báo cáo đồ án được biên soạn bằng \textbf{\LaTeX} trên
Overleaf theo đúng mẫu quy định.

% =====================================================================
\section{Kết chương}
\label{sec:ket-chuong-3}

Chương này đã trình bày các cơ sở lý thuyết --- ôn tập ngắt quãng và thuật
toán SM-2, khung trình độ CEFR, kiến trúc REST, ánh xạ đối tượng--quan hệ và
mô hình ngôn ngữ lớn --- cùng ngăn xếp công nghệ được lựa chọn cho từng tầng
của hệ thống, từ cơ sở dữ liệu, frontend, backend, cơ chế xác thực và phân
quyền, các dịch vụ trí tuệ nhân tạo đến bộ công cụ phát triển. Những lựa
chọn này trực tiếp phục vụ các quyết định thiết kế ở
Chương~\ref{chuong:phan-tich-thiet-ke} và là tiền đề cho quá trình hiện
thực hoá và kiểm thử được trình bày ở chương tiếp theo.

%  vào main.tex (sau Chương 2).
%
%  GHI CHÚ VỀ HÌNH VẼ: dùng macro \figph{<tên-file>}{<tỉ-lệ-rộng>}{<chú-thích>}{<nhãn>}
%  (đã khai báo ở Chương 2, khai báo lại bằng \providecommand để an toàn nếu
%  biên dịch độc lập). Macro tự chèn ảnh thật nếu tệp tồn tại trong images/,
%  ngược lại hiển thị khung giữ chỗ để báo cáo vẫn biên dịch được.
%
%  KHÔNG hardcode số chương/mục/bảng/hình/công thức — luôn dùng \label + \ref / \cite.
% =====================================================================

% --- Macro tiện ích (idempotent — không định nghĩa lại nếu đã có) ------
\providecommand{\code}[1]{\texttt{\detokenize{#1}}} % in định danh snake_case an toàn
\providecommand{\figph}[4]{%
  \begin{figure}[htbp]
    \centering
    \IfFileExists{images/#1}%
      {\includegraphics[width=#2\textwidth,height=0.82\textheight,keepaspectratio]{#1}}%
      {\setlength{\fboxsep}{10pt}\fbox{\begin{minipage}[c][3.6cm][c]{#2\textwidth}\centering\itshape Sơ đồ minh hoạ --- chèn tệp ảnh \code{images/#1} vào thư mục \code{images/}.\end{minipage}}}%
    \caption{#3}%
    \label{#4}%
  \end{figure}}

\chapter{CƠ SỞ LÝ THUYẾT VÀ CÔNG NGHỆ SỬ DỤNG}
\label{chuong:co-so-ly-thuyet}

Chương~\ref{chuong:phan-tich-thiet-ke} đã trình bày kết quả phân tích và
thiết kế hệ thống. Chương này hệ thống hoá những \textit{cơ sở lý thuyết}
và \textit{công nghệ} làm nền tảng cho các quyết định thiết kế và hiện thực
hoá đó, qua đó lý giải vì sao mỗi công nghệ được lựa chọn. Trước hết,
mục~\ref{sec:co-so-ly-thuyet} trình bày các cơ sở lý thuyết cốt lõi: ôn tập
ngắt quãng và thuật toán SM-2, khung trình độ CEFR, kiến trúc REST, ánh xạ
đối tượng--quan hệ và mô hình ngôn ngữ lớn. Tiếp theo,
mục~\ref{sec:tong-quan-cong-nghe} tổng quan ngăn xếp công nghệ và
mục~\ref{sec:kien-truc-tong-quan} phân tích các kiểu kiến trúc được áp dụng.
Các mục~\ref{sec:cong-nghe-csdl}--\ref{sec:cong-cu-phat-trien} đi sâu lần
lượt vào công nghệ cơ sở dữ liệu, frontend, backend, cơ chế xác thực và
phân quyền, công nghệ trí tuệ nhân tạo và bộ công cụ phát triển.

% =====================================================================
\section{Cơ sở lý thuyết}
\label{sec:co-so-ly-thuyet}

\subsection{Ôn tập ngắt quãng và đường cong lãng quên}
\label{subsec:srs-theory}

Cơ sở khoa học của động cơ học trong hệ thống là hiện tượng
\textit{đường cong lãng quên} (forgetting curve) do Ebbinghaus phát hiện:
khả năng ghi nhớ một thông tin mới suy giảm theo hàm mũ theo thời gian nếu
không được củng cố \cite{ebbinghaus1913memory}. Một cách mô hình hoá đơn
giản, mức ghi nhớ \( R \) còn lại sau khoảng thời gian \( t \) kể từ lần
học gần nhất được biểu diễn qua công thức~\eqref{eq:forgetting}:

\begin{equation}
  \label{eq:forgetting}
  R(t) = e^{-\,t/S},
\end{equation}

trong đó \( S \) là \textit{độ bền trí nhớ} (memory stability) --- càng ôn
tập đúng lúc, \( S \) càng lớn và thông tin được giữ lâu hơn. Hệ quả thực
tiễn là: thay vì học nhồi nhét, người học nên ôn lại mỗi từ tại những
\textit{thời điểm đến hạn} được giãn dần. Các nghiên cứu về hiệu ứng giãn
cách (spacing effect) khẳng định ôn tập ngắt quãng cải thiện rõ rệt khả
năng ghi nhớ dài hạn so với học dồn \cite{cepeda2006distributed}. Đây là
nền tảng lý thuyết cho \textit{hệ thống ôn tập ngắt quãng} (Spaced
Repetition System --- SRS) mà hệ thống hiện thực hoá.

\subsection{Thuật toán lập lịch ôn tập SM-2}
\label{subsec:sm2-theory}

Để quyết định \textit{khi nào} ôn lại mỗi từ, hệ thống áp dụng thuật toán
\textbf{SM-2} có nguồn gốc từ dự án SuperMemo \cite{wozniak1990optimization}.
Mỗi cặp (người dùng, từ) được mô tả bằng ba đại lượng: hệ số dễ nhớ
\textit{easiness factor} \( EF \) (khởi tạo \( 2{,}5 \)), số lần ôn đúng
liên tiếp \( n \) và khoảng cách ôn tập \( I \) (tính theo ngày). Sau mỗi
lượt ôn, người học (ở đây là kết quả trả lời trong phiên) được quy về một
\textit{điểm chất lượng} \( q \in \{0,1,\dots,5\} \).

Hệ số \( EF \) được cập nhật theo công thức~\eqref{eq:sm2-ef}, có chặn dưới
\( 1{,}3 \) để khoảng ôn tập không co lại quá nhanh:

\begin{equation}
  \label{eq:sm2-ef}
  EF' = EF + \big(0{,}1 - (5-q)\,(0{,}08 + (5-q)\cdot 0{,}02)\big),
  \qquad EF' \geq 1{,}3 .
\end{equation}

Khoảng cách ôn tập kế tiếp được tính theo công thức~\eqref{eq:sm2-interval}:

\begin{equation}
  \label{eq:sm2-interval}
  I(n) =
  \begin{cases}
    1, & n = 1,\\[2pt]
    6, & n = 2,\\[2pt]
    \big\lceil I(n-1)\times EF' \big\rceil, & n > 2 .
  \end{cases}
\end{equation}

Nếu người học trả lời sai (\( q < 3 \)), bộ đếm \( n \) được đặt lại về 0 và
từ phải học lại từ đầu; ngược lại \( n \) tăng và khoảng ôn tập giãn dần
theo \( EF' \). Trong hệ thống, SM-2 được \textbf{mở rộng bằng các bước học}
(learning steps): một từ mới phải vượt qua một chuỗi câu hỏi từ dễ đến khó
trong phiên trước khi được giao cho SM-2 lập lịch dài hạn; điểm chất lượng
\( q \) được suy ra từ độ chính xác ở bước cuối của mỗi từ. Trạng thái thẻ
nhờ đó chuyển dịch \code{new} \( \to \) \code{learning} \( \to \)
\code{review} \( \to \) \code{mastered}, phù hợp với biểu đồ hoạt động đã
trình bày ở Chương~\ref{chuong:phan-tich-thiet-ke}.

\subsection{Khung tham chiếu trình độ ngôn ngữ CEFR}
\label{subsec:cefr-theory}

Hệ thống chuẩn hoá độ khó của từ vựng và trình độ người học theo
\textbf{Khung tham chiếu trình độ ngôn ngữ chung châu Âu} (Common European
Framework of Reference for Languages --- CEFR) do Hội đồng châu Âu ban hành
\cite{coe2001cefr}. CEFR chia năng lực ngôn ngữ thành sáu bậc từ
\( A1, A2 \) (sơ cấp), \( B1, B2 \) (trung cấp) đến \( C1, C2 \) (cao cấp).
Mỗi từ trong kho được gắn một bậc CEFR và mỗi người học khai báo trình độ
của mình; nhờ đó hệ thống lựa chọn từ và gợi ý bộ thẻ phù hợp với năng lực
hiện tại, đồng thời mô hình ngôn ngữ chấm điểm bậc CEFR mà một câu luyện
tập thể hiện.

\subsection{Kiến trúc REST và dịch vụ web}
\label{subsec:rest-theory}

Giao diện lập trình ứng dụng (API) của hệ thống tuân theo phong cách kiến
trúc \textbf{REST} (Representational State Transfer) do Fielding đề xuất
\cite{fielding2000rest}. Các ràng buộc cốt lõi của REST mà hệ thống áp dụng
gồm: mô hình \textit{client--server} tách bạch máy khách và máy chủ; tính
\textit{phi trạng thái} (stateless) --- mỗi yêu cầu tự mang đủ thông tin xác
thực (token), máy chủ không lưu trạng thái phiên; định danh tài nguyên bằng
\textit{URI} và thao tác qua các phương thức HTTP chuẩn (GET, POST, PATCH,
DELETE) với ngữ nghĩa thống nhất; và trao đổi dữ liệu dưới dạng
\textbf{JSON}. Tính phi trạng thái là điều kiện để mở rộng tầng API theo
chiều ngang (đáp ứng yêu cầu NFR-03 trong
Bảng~\ref{tab:nfr}). API còn được \textit{phiên bản hoá theo URI}
(tiền tố \code{/v1}) nhằm tiến hoá hợp đồng mà không phá vỡ máy khách cũ.

\subsection{Ánh xạ đối tượng--quan hệ (ORM)}
\label{subsec:orm-theory}

Để thu hẹp khoảng cách giữa mô hình hướng đối tượng trong mã nguồn và mô
hình quan hệ trong cơ sở dữ liệu, hệ thống dùng kỹ thuật \textbf{ánh xạ đối
tượng--quan hệ} (Object--Relational Mapping --- ORM). Mỗi bảng được khai
báo dưới dạng một lớp \textit{thực thể} (entity), mỗi dòng tương ứng một đối
tượng và các quan hệ khoá ngoại được biểu diễn bằng liên kết giữa các thực
thể. ORM giúp lập trình viên thao tác dữ liệu bằng mã có kiểu (typed) thay
vì viết câu lệnh SQL thủ công, đồng thời quản lý tiến hoá lược đồ bằng cơ
chế \textit{migration} có thể phiên bản hoá. Đánh đổi của ORM là một lớp
trừu tượng phụ và nguy cơ truy vấn kém hiệu quả nếu lạm dụng; hệ thống khắc
phục bằng cách đánh chỉ mục các truy vấn nóng và truy vấn tường minh khi cần.

\subsection{Mô hình ngôn ngữ lớn và sinh nội dung}
\label{subsec:llm-theory}

Các chức năng hỗ trợ bằng AI (làm giàu dữ liệu từ điển, chấm điểm câu luyện
tập) dựa trên \textbf{mô hình ngôn ngữ lớn} (Large Language Model --- LLM).
LLM là mạng nơ-ron kiến trúc \textit{Transformer} được huấn luyện trên khối
dữ liệu văn bản khổng lồ để dự đoán token kế tiếp, nhờ đó có khả năng sinh
văn bản mạch lạc, dịch và đánh giá ngôn ngữ theo chỉ dẫn (prompt). Hệ thống
sử dụng họ mô hình mở \textbf{Gemma} của Google DeepMind
\cite{gemma2025technical}. Trong nghiệp vụ, LLM được giao vai trò có kiểm
soát: sinh nghĩa, ví dụ, bản dịch theo định dạng yêu cầu và chấm điểm câu
theo \textit{rubric}; kết quả do AI sinh ra đối với kho hệ thống luôn phải
qua bước \textit{duyệt} của quản trị viên trước khi xuất bản, nhằm kiểm soát
chất lượng và hạn chế lỗi "ảo giác" (hallucination) của mô hình.

% =====================================================================
\section{Tổng quan công nghệ sử dụng}
\label{sec:tong-quan-cong-nghe}

Trên cơ sở lý thuyết ở mục~\ref{sec:co-so-ly-thuyet}, hệ thống lựa chọn ngăn
xếp công nghệ được tổ chức theo nhiều tầng. Bảng~\ref{tab:tech-overview}
tóm tắt công nghệ chủ chốt của từng tầng cùng vai trò của nó;
Hình~\ref{fig:tech-stack} minh hoạ ngăn xếp này một cách trực quan. Các mục
tiếp theo phân tích chi tiết từng nhóm công nghệ.

\begin{table}[h]
  \centering
  \small
  \renewcommand{\arraystretch}{1.3}
  \begin{tabular}{|p{3.0cm}|p{5.0cm}|p{6.4cm}|}
    \hline
    \textbf{Tầng} & \textbf{Công nghệ chính} & \textbf{Vai trò} \\
    \hline
    Frontend & Next.js, React, TypeScript, Tailwind CSS, shadcn/ui & Ứng dụng khách tiêu thụ API, dựng giao diện người dùng. \\
    \hline
    Backend & NestJS 11, TypeScript, TypeORM & Cung cấp API REST, chứa nghiệp vụ và truy cập dữ liệu. \\
    \hline
    Cơ sở dữ liệu & PostgreSQL & Lưu trữ dữ liệu quan hệ có ràng buộc toàn vẹn. \\
    \hline
    Hàng đợi \& nền & Redis, BullMQ & Hàng đợi công việc và worker xử lý bất đồng bộ. \\
    \hline
    Xác thực & JWT, Passport, bcrypt, OAuth/OIDC & Xác thực, phân quyền và bảo vệ tài khoản. \\
    \hline
    Trí tuệ nhân tạo & Gemma~3 (Google AI Studio), Edge TTS, vi dịch vụ chấm phát âm & Sinh nội dung, sinh âm thanh, chấm điểm luyện tập. \\
    \hline
    Hạ tầng phụ trợ & Cloudinary, SMTP, Docker & Lưu trữ media, gửi email, đóng gói triển khai. \\
    \hline
  \end{tabular}
  \caption{Tổng quan ngăn xếp công nghệ của hệ thống}
  \label{tab:tech-overview}
\end{table}

% --- PlantUML gợi ý cho Hình ngăn xếp công nghệ ----------------------
% @startuml
% rectangle "FRONTEND\nNext.js + React + TypeScript\nTailwind CSS + shadcn/ui" as FE
% rectangle "BACKEND (API)\nNestJS 11 + TypeScript\nController - Service - Repository (TypeORM)" as BE
% rectangle "XỬ LÝ NỀN\nWorker (BullMQ) + Redis" as WK
% database "PostgreSQL" as DB
% cloud "DỊCH VỤ AI / NGOÀI\nGemma 3, Edge TTS, Cloudinary,\nDV chấm phát âm, OAuth, SMTP" as EXT
% FE --> BE : HTTPS / JSON (REST)
% BE --> DB
% BE --> WK : enqueue
% WK --> DB
% WK --> EXT
% BE --> EXT
% @enduml
\figph{tech-stack.png}{0.8}{Ngăn xếp công nghệ của hệ thống theo tầng}{fig:tech-stack}

% =====================================================================
\section{Tổng quan kiến trúc hệ thống}
\label{sec:kien-truc-tong-quan}

Mục này điểm lại các \textit{kiểu kiến trúc} (architectural styles) làm nền
cho thiết kế đã trình bày chi tiết ở mục kiến trúc của
Chương~\ref{chuong:phan-tich-thiet-ke} (xem Hình~\ref{fig:kien-truc}).

\begin{itemize}
  \item \textbf{Monolith mô-đun hoá (modular monolith):} toàn bộ backend là
        một tiến trình triển khai duy nhất nhưng được chia thành các mô-đun
        nghiệp vụ tách bạch (xem Bảng~\ref{tab:modules}). So với kiến trúc
        vi dịch vụ (microservices), cách này đơn giản hơn khi triển khai và
        vận hành ở quy mô đồ án, đồng thời vẫn giữ ranh giới mô-đun rõ ràng
        để có thể tách dịch vụ về sau.
  \item \textbf{Kiến trúc phân tầng (layered):} trong mỗi mô-đun, mã nguồn
        chia theo các tầng \textit{Controller} (tiếp nhận HTTP, kiểm tra
        DTO) \( \to \) \textit{Service} (nghiệp vụ) \( \to \)
        \textit{Repository} (truy cập dữ liệu qua ORM). Việc tách tầng giúp
        kiểm thử và bảo trì độc lập từng mối quan tâm.
  \item \textbf{Client--server phi trạng thái (REST):} máy khách
        (Next.js) và máy chủ giao tiếp qua HTTP/JSON theo REST; máy chủ phi
        trạng thái nên có thể nhân bản để mở rộng theo chiều ngang.
  \item \textbf{Xử lý bất đồng bộ kiểu nhà sản xuất--người tiêu thụ
        (producer--consumer):} tiến trình API đóng vai trò nhà sản xuất, đẩy
        công việc nặng vào hàng đợi; tiến trình worker là người tiêu thụ,
        xử lý nền. Mẫu này giữ cho luồng yêu cầu chính luôn nhanh và ổn định.
\end{itemize}

Sự kết hợp bốn kiểu kiến trúc trên cho phép hệ thống vừa đơn giản khi triển
khai, vừa đáp ứng các yêu cầu phi chức năng về hiệu năng, khả năng mở rộng
và khả năng bảo trì trong Bảng~\ref{tab:nfr}.

% =====================================================================
\section{Công nghệ cơ sở dữ liệu}
\label{sec:cong-nghe-csdl}

Hệ thống dùng hệ quản trị cơ sở dữ liệu quan hệ \textbf{PostgreSQL}
\cite{postgresql2024}. PostgreSQL được chọn vì là mã nguồn mở, trưởng thành,
tuân thủ chuẩn SQL và hỗ trợ tốt các đặc tính mà thiết kế dữ liệu
(mục thiết kế CSDL của Chương~\ref{chuong:phan-tich-thiet-ke}) yêu cầu:

\begin{itemize}
  \item \textbf{Giao dịch ACID:} bảo đảm tính nguyên tử cho các thao tác
        ghi nhiều bảng (ví dụ tạo phiên học, cập nhật tiến độ và ghi nhật
        ký), đáp ứng yêu cầu độ tin cậy.
  \item \textbf{Khoá chính \code{uuid}:} dùng kiểu định danh duy nhất toàn
        cục cho mọi bảng, thuận lợi cho phân tán và tránh lộ số lượng bản ghi.
  \item \textbf{Kiểu \code{timestamptz}:} lưu mốc thời gian kèm múi giờ,
        bảo đảm nhất quán cho thống kê chuỗi ngày học và lịch ôn tập.
  \item \textbf{Ràng buộc toàn vẹn:} khoá ngoại, ràng buộc duy nhất và quy
        tắc xoá lan truyền (\code{CASCADE} / \code{SET NULL}) giữ dữ liệu
        nhất quán giữa từ vựng, bộ thẻ và tiến độ.
  \item \textbf{Chỉ mục:} đánh chỉ mục các cột truy vấn nóng (lemma, trạng
        thái ôn tập, thời điểm đến hạn) để bảo đảm hiệu năng khi dữ liệu lớn.
\end{itemize}

Bên cạnh PostgreSQL, hệ thống dùng \textbf{Redis} làm kho khoá--giá trị
trong bộ nhớ, đóng vai trò \textit{nơi lưu trữ hàng đợi} cho BullMQ
\cite{redis2024}. Việc truy cập cơ sở dữ liệu từ phía backend được thực hiện
qua \textbf{TypeORM} --- một thư viện ORM cho TypeScript/Node.js
\cite{typeorm2024} --- sẽ được trình bày ở mục~\ref{sec:cong-nghe-backend}.
Lược đồ cơ sở dữ liệu được quản lý bằng các \textit{migration} đặt trong thư
mục mã nguồn và chạy qua công cụ dòng lệnh của TypeORM, giúp tiến hoá lược
đồ một cách có kiểm soát và tái lập được giữa các môi trường.

% =====================================================================
\section{Công nghệ Frontend}
\label{sec:cong-nghe-frontend}

Ứng dụng khách (frontend) là một dự án độc lập, đóng vai trò bên tiêu thụ
API REST của backend. Frontend được xây dựng bằng \textbf{Next.js 16} theo
mô hình \textit{App Router} --- framework React hỗ trợ kết xuất phía máy chủ
(Server-Side Rendering --- SSR), \textit{thành phần phía máy chủ} (Server
Components) và \textit{hành động phía máy chủ} (Server Actions)
\cite{nextjs2024} --- trên nền thư viện giao diện \textbf{React~19}
\cite{react2024} và ngôn ngữ \textbf{TypeScript} (chế độ \textit{strict}).
Việc dùng chung TypeScript ở cả frontend và backend giúp thống nhất kiểu dữ
liệu của hợp đồng API giữa hai phía. Phần tạo kiểu giao diện dùng
\textbf{Tailwind CSS v4} --- framework CSS theo hướng tiện ích (utility-first),
cấu hình bằng CSS \cite{tailwind2024} --- kết hợp bộ thành phần
\textbf{shadcn/ui} dựng trên các primitive khả truy cập của \textbf{Base UI}.
Ràng buộc và kiểm tra dữ liệu biểu mẫu dùng \code{zod} \cite{zod2024}, được
dùng lại giữa kiểm tra phía máy khách và Server Action. Bảng~\ref{tab:frontend-libs}
liệt kê các thư viện frontend chính.

\begin{table}[h]
  \centering
  \small
  \renewcommand{\arraystretch}{1.3}
  \begin{tabular}{|p{4.2cm}|p{10.4cm}|}
    \hline
    \textbf{Công nghệ / thư viện} & \textbf{Vai trò} \\
    \hline
    Next.js 16 (App Router), React 19, TypeScript & Khung ứng dụng theo App Router; Server Components/Server Actions, định tuyến và kết xuất trang. \\
    \hline
    Tailwind CSS v4, shadcn/ui, Base UI & Hệ thống tạo kiểu tiện ích và bộ thành phần giao diện khả truy cập, tái sử dụng. \\
    \hline
    \code{zod} & Khai báo và kiểm tra ràng buộc dữ liệu (biểu mẫu, phản hồi API) phía máy khách. \\
    \hline
    \code{lucide-react}, \code{sonner} & Bộ biểu tượng và thông báo nổi (toast) cho phản hồi thao tác. \\
    \hline
    \code{next-themes}, \code{input-otp} & Hỗ trợ chế độ sáng/tối và ô nhập mã OTP (xác minh email). \\
    \hline
  \end{tabular}
  \caption{Các công nghệ và thư viện frontend chính}
  \label{tab:frontend-libs}
\end{table}

Mục này chỉ giới thiệu \textit{ngăn xếp công nghệ} frontend; quá trình hiện
thực hoá tầng giao diện (cấu trúc App Router, tầng truy cập API phía máy chủ,
quản lý phiên và động cơ phiên học phía máy khách) được trình bày ở
mục~\ref{sec:hien-thuc-frontend}, Chương~\ref{chuong:hien-thuc-hoa}. Trong mọi
trường hợp, frontend giao tiếp với backend hoàn toàn qua các API đã đặc tả,
không truy cập trực tiếp cơ sở dữ liệu.

% =====================================================================
\section{Công nghệ Backend}
\label{sec:cong-nghe-backend}

Backend được xây dựng bằng \textbf{NestJS}, một framework Node.js cho
TypeScript theo hướng mô-đun, dựa trên các nguyên lý \textit{tiêm phụ thuộc}
(dependency injection) và lập trình hướng khía cạnh \cite{nestjs2024}.
NestJS phù hợp với kiến trúc phân tầng đã nêu ở
mục~\ref{sec:kien-truc-tong-quan}: mỗi mô-đun đóng gói controller, service
và provider riêng; các \textit{guard}, \textit{pipe} và
\textit{interceptor} xử lý những mối quan tâm xuyên suốt (xác thực, kiểm tra
đầu vào, biến đổi phản hồi) một cách khai báo. Toàn bộ mã nguồn viết bằng
\textbf{TypeScript} ở chế độ \textit{strict}, giúp phát hiện lỗi kiểu ngay
khi biên dịch. Bảng~\ref{tab:backend-libs} liệt kê các thành phần backend
chính.

\begin{table}[h]
  \centering
  \small
  \renewcommand{\arraystretch}{1.3}
  \begin{tabular}{|p{4.6cm}|p{10.0cm}|}
    \hline
    \textbf{Công nghệ / thư viện} & \textbf{Vai trò} \\
    \hline
    NestJS 11 & Framework backend: tổ chức mô-đun, định tuyến, tiêm phụ thuộc, phiên bản hoá API theo URI. \\
    \hline
    TypeORM & Ánh xạ đối tượng--quan hệ tới PostgreSQL; quản lý thực thể và migration. \\
    \hline
    \code{class-validator}, \code{class-transformer} & Kiểm tra ràng buộc DTO và biến đổi dữ liệu vào/ra qua \code{ValidationPipe} toàn cục (loại bỏ trường lạ). \\
    \hline
    BullMQ + Redis & Hàng đợi công việc và worker xử lý nền (làm giàu từ vựng, sinh âm thanh, chấm điểm luyện tập). \\
    \hline
    \code{@nestjs/throttler} & Giới hạn tần suất (rate limiting) cho các điểm cuối nhạy cảm như đăng nhập. \\
    \hline
    \code{nodemailer} & Gửi email (mã xác minh) qua SMTP. \\
    \hline
    \code{pdf-parse}, \code{xlsx} & Trích xuất văn bản từ tệp PDF và bảng tính phục vụ nhập từ vựng hàng loạt. \\
    \hline
  \end{tabular}
  \caption{Các công nghệ và thư viện backend chính}
  \label{tab:backend-libs}
\end{table}

Phần xử lý bất đồng bộ được hiện thực bằng \textbf{BullMQ} --- thư viện hàng
đợi công việc chạy trên Redis \cite{bullmq2024}. Tiến trình API và tiến
trình worker chia sẻ chung cơ sở dữ liệu và hàng đợi nhưng được khởi chạy
độc lập (lệnh \code{start} và \code{start:worker}), phản ánh đúng mô hình
producer--consumer ở mục~\ref{sec:kien-truc-tong-quan}.

% =====================================================================
\section{Cơ chế xác thực và phân quyền}
\label{sec:xac-thuc-phan-quyen}

Phân hệ tài khoản áp dụng nhiều lớp kỹ thuật bảo mật, đáp ứng yêu cầu
NFR-01 trong Bảng~\ref{tab:nfr}.

\subsection{Băm mật khẩu}
\label{subsec:bcrypt}

Mật khẩu người dùng không bao giờ lưu ở dạng văn bản thuần mà được băm bằng
\textbf{bcrypt} --- hàm băm mật khẩu thích nghi, có \textit{muối} (salt) và
hệ số chi phí điều chỉnh được để chống tấn công vét cạn
\cite{provos1999bcrypt}. Khi đăng nhập, hệ thống so khớp mật khẩu nhập vào
với bản băm đã lưu thay vì giải mã.

\subsection{Xác thực bằng JSON Web Token}
\label{subsec:jwt}

Sau khi xác thực thành công, hệ thống cấp một cặp token theo chuẩn
\textbf{JSON Web Token (JWT)} \cite{rfc7519}: một \textit{access token} sống
ngắn để gọi API và một \textit{refresh token} sống dài để làm mới phiên.
JWT là chuỗi tự chứa, được ký số nên máy chủ có thể xác minh tính toàn vẹn
mà không cần tra cứu trạng thái phiên --- đúng tinh thần phi trạng thái của
REST (mục~\ref{subsec:rest-theory}). Thư viện \textbf{Passport} (qua
\code{passport-jwt}) đảm nhiệm việc trích và xác minh token trong mỗi yêu
cầu. Để có thể \textit{thu hồi}, refresh token được lưu dưới dạng băm trong
bảng \code{refresh_tokens} và bị vô hiệu khi người dùng đăng xuất.

\subsection{Phân quyền theo vai trò và quyền sở hữu}
\label{subsec:rbac}

Việc cấp quyền dựa trên hai cơ chế bổ trợ:

\begin{itemize}
  \item \textbf{Theo vai trò (role-based):} mỗi tài khoản mang vai trò
        \code{user} hoặc \code{admin}; một \textit{guard} kiểm tra vai trò
        để chặn các thao tác quản trị đối với người dùng thường.
  \item \textbf{Theo quyền sở hữu (ownership-based):} với tài nguyên cá
        nhân (từ vựng riêng, bộ thẻ, tiến độ), hệ thống kiểm tra người gọi
        có đúng là chủ sở hữu hay không trước khi cho phép truy cập.
\end{itemize}

\subsection{Đăng nhập qua nhà cung cấp bên thứ ba (OAuth/OIDC)}
\label{subsec:oauth}

Ngoài email/mật khẩu, hệ thống cho phép đăng nhập qua \textbf{Google},
\textbf{Apple} và \textbf{GitHub} bằng giao thức \textbf{OAuth~2.0} và lớp
nhận dạng \textbf{OpenID Connect} (OIDC) \cite{rfc6749}. Token danh tính do
nhà cung cấp cấp được hệ thống xác minh, sau đó liên kết với một bản ghi
trong bảng \code{user_identities}; nhờ vậy một người dùng có thể đăng nhập
bằng nhiều phương thức mà vẫn chung một tài khoản.

\subsection{Chấm điểm phiên học phi trạng thái bằng HMAC}
\label{subsec:hmac}

Một điểm đặc thù về bảo mật là cơ chế \textbf{chấm điểm phi trạng thái} cho
phiên học. Thay vì lưu trạng thái phiên trên máy chủ, mỗi câu hỏi khi sinh
ra được \textit{ký} bằng \textbf{HMAC} (Hash-based Message Authentication
Code) \cite{rfc2104} dựa trên một khoá bí mật của máy chủ. Khi người học nộp
đáp án, máy chủ xác minh lại chữ ký để bảo đảm câu hỏi không bị giả mạo hay
hết hạn, rồi mới suy lại đáp án đúng và chấm điểm. Cách làm này vừa giữ tính
phi trạng thái (thuận lợi cho mở rộng ngang) vừa chống gian lận.

% =====================================================================
\section{Công nghệ trí tuệ nhân tạo}
\label{sec:cong-nghe-ai}

Các chức năng thông minh của hệ thống dựa trên một số dịch vụ AI được tích
hợp qua tầng worker xử lý nền. Bảng~\ref{tab:ai-services} tổng hợp các dịch
vụ này.

\begin{table}[h]
  \centering
  \small
  \renewcommand{\arraystretch}{1.3}
  \begin{tabular}{|p{3.6cm}|p{4.0cm}|p{6.8cm}|}
    \hline
    \textbf{Dịch vụ} & \textbf{Loại} & \textbf{Vai trò trong hệ thống} \\
    \hline
    Gemma~3 (Google AI Studio) & Mô hình ngôn ngữ lớn (LLM) & Sinh nghĩa, ví dụ, bản dịch; chấm điểm câu luyện viết theo rubric kèm bậc CEFR. \\
    \hline
    Microsoft Edge TTS & Tổng hợp giọng nói (TTS) & Sinh âm thanh phát âm tự động cho từ vựng. \\
    \hline
    Vi dịch vụ chấm phát âm & Mô hình chấm điểm âm vị & Nhận bản ghi âm, trả điểm tổng và điểm từng phoneme. \\
    \hline
  \end{tabular}
  \caption{Các dịch vụ trí tuệ nhân tạo được tích hợp}
  \label{tab:ai-services}
\end{table}

\subsection{Sinh nội dung và chấm điểm bằng Gemma~3}
\label{subsec:gemma}

Hệ thống gọi mô hình \textbf{Gemma~3} qua nền tảng Google AI Studio
\cite{gemma2025technical}. Mô hình được sử dụng cho hai nhóm tác vụ:
(i)~\textit{làm giàu dữ liệu từ điển} --- từ một lemma, sinh các nghĩa, câu
ví dụ và bản dịch theo định dạng có cấu trúc; và (ii)~\textit{chấm điểm câu
luyện tập} --- đánh giá câu người học đặt theo thang điểm 0--100 kèm bậc
CEFR và nhận xét. Mọi lời gọi LLM đều là tác vụ nặng nên được đẩy sang
worker xử lý bất đồng bộ và bị giới hạn theo hạn mức ngày để bảo vệ khoá API
dùng chung.

\subsection{Tổng hợp giọng nói và chấm điểm phát âm}
\label{subsec:tts-pron}

Âm thanh phát âm của từ vựng được sinh tự động bằng \textbf{Microsoft Edge
TTS} (qua thư viện \code{msedge-tts}) rồi tải lên kho media. Chức năng luyện
phát âm được phục vụ bởi một \textbf{vi dịch vụ} riêng viết bằng Python
(FastAPI): backend đóng vai trò proxy mỏng, chuyển tiếp bản ghi âm của người
học tới vi dịch vụ để nhận điểm tổng và điểm từng âm vị, sau đó lưu lại lịch
sử luyện tập. Việc tách phần chấm phát âm thành dịch vụ độc lập giúp cô lập
các phụ thuộc nặng về xử lý tín hiệu/học máy khỏi tiến trình backend chính.

% =====================================================================
\section{Công cụ phát triển}
\label{sec:cong-cu-phat-trien}

Quá trình phát triển sử dụng bộ công cụ trong Bảng~\ref{tab:dev-tools} nhằm
bảo đảm chất lượng mã nguồn, khả năng tái lập môi trường và năng suất làm
việc.

\begin{table}[h]
  \centering
  \small
  \renewcommand{\arraystretch}{1.3}
  \begin{tabular}{|p{4.0cm}|p{10.6cm}|}
    \hline
    \textbf{Công cụ} & \textbf{Mục đích} \\
    \hline
    Git \& GitHub & Quản lý phiên bản mã nguồn và cộng tác qua pull request. \\
    \hline
    Docker \& Docker Compose & Đóng gói và khởi chạy nhất quán PostgreSQL, Redis, pgAdmin phục vụ phát triển và triển khai \cite{docker2024}. \\
    \hline
    NestJS CLI & Sinh khung mô-đun/controller/service và biên dịch, chạy ứng dụng. \\
    \hline
    TypeORM CLI & Sinh, chạy và hoàn tác migration của lược đồ cơ sở dữ liệu. \\
    \hline
    ESLint \& Prettier & Phân tích tĩnh và định dạng mã TypeScript theo quy ước thống nhất. \\
    \hline
    Jest \& Supertest & Kiểm thử đơn vị và kiểm thử tích hợp (end-to-end) cho API. \\
    \hline
    pgAdmin & Quản trị và truy vấn trực quan cơ sở dữ liệu PostgreSQL. \\
    \hline
    Overleaf (\LaTeX) & Biên soạn báo cáo đồ án theo mẫu của nhà trường. \\
    \hline
  \end{tabular}
  \caption{Các công cụ phát triển sử dụng trong đề tài}
  \label{tab:dev-tools}
\end{table}

Mã nguồn được quản lý bằng \textbf{Git} và lưu trên GitHub; mỗi tính năng
hoặc sửa lỗi được phát triển trên một nhánh riêng rồi hợp nhất qua pull
request. Môi trường phụ thuộc (PostgreSQL, Redis, pgAdmin) được mô tả bằng
\textbf{Docker Compose} \cite{docker2024}, giúp tái lập nhanh và đồng nhất
giữa các máy. Chất lượng mã được kiểm soát bằng kiểm tra kiểu của TypeScript,
\textbf{ESLint} và \textbf{Prettier}; phần kiểm thử dùng \textbf{Jest} cho
kiểm thử đơn vị và \textbf{Supertest} cho kiểm thử tích hợp các điểm cuối
API. Cuối cùng, báo cáo đồ án được biên soạn bằng \textbf{\LaTeX} trên
Overleaf theo đúng mẫu quy định.

% =====================================================================
\section{Kết chương}
\label{sec:ket-chuong-3}

Chương này đã trình bày các cơ sở lý thuyết --- ôn tập ngắt quãng và thuật
toán SM-2, khung trình độ CEFR, kiến trúc REST, ánh xạ đối tượng--quan hệ và
mô hình ngôn ngữ lớn --- cùng ngăn xếp công nghệ được lựa chọn cho từng tầng
của hệ thống, từ cơ sở dữ liệu, frontend, backend, cơ chế xác thực và phân
quyền, các dịch vụ trí tuệ nhân tạo đến bộ công cụ phát triển. Những lựa
chọn này trực tiếp phục vụ các quyết định thiết kế ở
Chương~\ref{chuong:phan-tich-thiet-ke} và là tiền đề cho quá trình hiện
thực hoá và kiểm thử được trình bày ở chương tiếp theo.
